\chapter{多模态模型}
\label{chap:vision-multimodal-models}

同一张街景照片，可以用来寻找相似图片、回答“右侧停着什么车”，也可以作为生成另一种天气下街景的条件。若再给一段录音，系统还能比较画面与声音是否相关，把回答说出来，甚至生成一段带发动机声的短视频。这些功能都接受多种输入，内部需要的计算却不同。一个能比较图文相似度的模型，未必有生成句子的输出层；一个会说话的视觉助手，也未必能生成视频。我们要回答的问题是：文字、图像和声音怎样进入同一个模型，在哪里发生联系，又怎样恢复成分数、语言或媒体？

\ref{chap:vision-image-video-models}章已经说明图像和视频怎样切块、编码，以及图像与视频生成主干怎样工作。本章在这些视觉表示之上加入文字和声音。先沿“编码—连接—解码”的计算流程固定共同的记号和接口；随后按输出类型讨论模型：只需比较时，用可比较表示完成理解；需要说出答案时，让语言模型读取视觉信息；需要创作媒体时，把文字或图像条件接入图像、视频和音频生成器。理解和生成各自成立后，又出现了能否把二者结合的问题，我们按结合紧密程度比较三种做法。最后回到同一街景和音画请求，检查信息、模态贡献和资源三方面的使用约束。阅读本章需要卷五的Transformer与注意力、卷六的扩散模型和变分自编码器，以及卷八的对比学习；用到时会先简要说明所用结论。多模态对齐与融合的一般原理已在\ref{chap:vision-multimodal-alignment}章讲过，本章把那里的对比对齐、注意力读取和中间融合落实为具体模型，并对每个模型写出层数、宽度、单元数等论文公开的结构配置；论文或报告没有公开的数字，表中注明未给出，不作推测。

以下沿用同一个教学场景：街景左侧有红色自行车，右侧有红色汽车，录音中可能出现发动机声；视频生成的例子沿用后文任务章中玩具狗走过桌面的片段。图片只提供静态外观，声音是否来自右车还需额外依据。这样，表示匹配、语言响应和媒体生成能够在相同对象上比较，而不会把生成的内容当作实拍证据。本章所用的小向量和概率均为教学构造，不是所介绍模型的测量结果。

\section{多模态模型的计算流程}
\label{sec:vision-models-interfaces}

多模态模型要完成的计算可以拆成三个先后环节。第一，把每种输入转成网络能处理的数值单元，并保留它们来自哪里、对应哪个位置或时刻；第二，让不同来源的单元发生联系，使一种模态的内容能够影响另一种模态的表示；第三，把联系后的内部表示变成任务需要的结果，例如分数、文字或媒体。三个环节各有难点：编码前的裁剪和采样已经限定了可用证据，连接方式决定哪些表示能预先计算、哪些必须随请求重算，输出端则需要与预测对象相容的解码器。后面各节的模型在这三个环节上作出不同选择，本节先把共同的记号和可选方案建立起来。

连接图像和文字之前，还需要区分原始观测、网络内部表示和最终结果。它们的数值形式可能相似，却承担不同职责。例如，一张照片可以形成许多空间位置上的特征，也可以被汇总为一个整图向量；后者便于匹配，却未必足以恢复前者包含的小字和边界。

\subsection{各模态输入的编码}
\label{sec:vision-models-encoding}

编码要解决的问题是：输入的形式各不相同，图像是像素网格，文字是字符串，声音是波形，网络却只处理定长向量组成的序列或网格。每种输入都需要一套固定的转换规则，把它切成计算单元、映射成向量，并记下每个单元对应的空间或时间范围。这种形式上的差别就是\ref{sec:vision-alignment-characteristics}节所说的表示异构：不同模态的取值类型、采样方式和坐标含义都不同，同一个“第$i$个单元”在图像中指一块空间区域，在声音中指一段时间。该节从数据特点出发讨论异构带来的对齐困难，本小节只交代各类模型实际采用的编码规则。

图像和视频的编码在\ref{sec:vision-ivm-pipeline}节已经展开。这里只回顾本章用到的结论：图像$\bm X\in\mathbb R^{h\times w\times c}$按边长$p$的不重叠块切分，经线性映射和位置表示得到$\bm F\in\mathbb R^{n_{\mathrm v}\times d_{\mathrm v}}$，$n_{\mathrm v}=hw/p^2$随分辨率增加；视频还要按帧或时空块组织，单元可以覆盖数帧。视觉Transformer（ViT）和卷积编码器的结构与训练分别见\ref{sec:vision-ivm-vit}节和\ref{sec:vision-ivm-resnet}节。本章只把$\bm F$当作视觉编码器给出的接口，关心它怎样与文字和声音相接。

文字先被切成离散\term{词元}（token），即分词器按词表规则划出的字符或子词单元，再经嵌入表转为向量。若词表大小为$V$，第$j$个词元编号为$a_j\in\{1,\ldots,V\}$，嵌入表$\bm B\in\mathbb R^{V\times d_{\mathrm t}}$给出的输入就是第$a_j$行。分词器规定哪些字符串对应哪些编号，训练和使用必须用同一套，不能在使用时任意换表。

声音以波形采样$s_i$记录，采样率为$f_{\mathrm s}$时第$i$个采样点对应$i/f_{\mathrm s}$秒。直接处理波形的序列很长：16kHz采样的10秒录音就有16万个点。而且单个采样点几乎没有含义，发动机声和人声的区别体现在一段时间内哪些频率的振动更强。常见做法是先把波形转成\term{梅尔频谱}（mel-spectrogram），即按听觉频带汇总的逐段频谱。它分两步计算。

第一步是\term{短时傅里叶变换}（short-time Fourier transform, STFT），即把长波形切成重叠短段、对每段分别求频谱的变换。整段录音只求一次频谱，会把开头的说话声和结尾的发动机声混在一起；切成短段后，每段内的声音可以近似看作不变，于是能得到“哪一时刻有哪些频率”。设每段长$N$个采样点、相邻两段起点相隔$H$个采样点，$w_i$是在段两端平滑降到零的窗函数，用来减弱截断造成的频谱泄漏，第$m$段第$k$个频率的系数为
\begin{equation}
  S_{m,k}=\sum_{i=0}^{N-1} s_{mH+i}\,w_i\,
          e^{-\mathrm{j}2\pi ki/N},
  \qquad k=0,1,\ldots,N/2.
  \label{eq:vision-models-stft}
\end{equation}
对固定的$m$，这就是对加窗片段做长度$N$的离散傅里叶变换。离散傅里叶变换把长度$N$的序列分解为$N$个复指数分量，第$k$个分量对应频率$kf_{\mathrm s}/N$；实信号的频谱正负频率共轭对称，所以只需保留$k\leq N/2$。这些结论可参考\BookChapterRef{chap:signal-analysis}{卷一“信号”章}，那里只讲一维整段序列，分段加窗是本章补充的用法。频率分辨率$f_{\mathrm s}/N$与时间分辨率$N/f_{\mathrm s}$此消彼长：段越长，相邻频率分得越细，但一段内混入的时间变化也越多。

第二步按听觉尺度合并频率。人耳对低频的音高差更敏感，对高频较迟钝，所以把频率轴换到\term{梅尔刻度}（mel scale）上，一种常用换算为$\mathrm{mel}(f)=2595\log_{10}(1+f/700)$。在梅尔刻度上等间隔放置$B$个三角形滤波器，每个滤波器对相邻若干频率的能量$|S_{m,k}|^2$加权求和，再取对数，得到第$m$段第$b$个频带的值。低频处的滤波器窄、高频处的宽，于是频带数$B$远小于$N/2+1$。

用本节常见的参数复算一遍。$f_{\mathrm s}=16000$，每段25毫秒即$N=400$个点，步长10毫秒即$H=160$个点；频率分辨率为$16000/400=40$Hz，每段有$N/2+1=201$个频率系数，合并成$B=128$个梅尔频带。10秒录音不补边时的段数为$\lfloor(160000-400)/160\rfloor+1=998$，约1000段，所以梅尔频谱约为$1000\times128$的数组。它的两维分别描述时间和频率，和图像的横纵坐标不同：沿时间轴平移一段表示事件提前或推迟，沿频率轴平移则改变音高；频率轴上相距较远的两个格子还可能属于同一个声音的谐波。之后可以像图像一样切块编码，但每个单元仍应记录对应的时间范围，后文ImageBind、Qwen2.5-Omni和LTX-2的音频路径都从这样的数组出发。

不同模态进入同一序列时，还要保留来源和坐标。实际系统常把连续的视觉特征也称为“视觉词元”，这里的词元表示可进入序列模型的一个计算单元，不一定经过离散词表量化。它与文字词元的整数编号、媒体生成器使用的离散编码、原始像素值都不能直接互换；比较词元数时，必须说明每个单元对应多少空间或时间范围。视频单纯把第1帧、第2帧写入序列，无法区分两帧相隔0.1秒还是10秒，因此还需记录真实时间戳；帧率与时间采样怎样限定可见的运动，见\ref{sec:vision-videofound-frame-rate}节。声音与画面要在时间上对照时，二者的单元间隔通常不同：上例的梅尔频谱每10毫秒一段，而每秒25帧的视频每40毫秒一帧，只有换算到共同的秒数，才能说某一段声音与某一帧同时发生。

坐标记录也决定了结果能否回到原始输入。设街景原图宽1200像素、高800像素，裁出左上角为$(400,200)$、宽高为$400\times400$的区域，再调整到$200\times200$作为局部输入。以连续边界坐标计，局部结果$(x',y')=(50,80)$对应原图
\[
  (x,y)=(400+2x',\,200+2y')=(500,360)\text{。}
\]
因此，输入记录至少应保留原图编号、裁剪边界和缩放关系。若视频只选取0、2、4秒三帧，还应保留真实时间戳，而不是让下游把三帧当成连续拍摄的相邻帧；同一段视频改抽1、3、5秒三帧，帧数相同，却可能包含不同的动作。裁剪与插值也应区分：从原始高分辨率照片重新裁取车窗，可能增加相对于缩小整图的可见细节；把已经缩小的车窗插值放大，仅改变数组尺寸，不增加原始像素证据。编码前的这些选择已经限定了模型可能看到的证据。

\subsection{跨模态信息的连接}
\label{sec:vision-models-connection}

编码后，图像、文字和声音各自成为一组向量，但这些向量还没有共同含义：图像编码器输出的第5维与文字编码器输出的第5维没有任何对应关系。连接要解决的问题是让一种模态的内容能够影响另一种模态的计算。\term{视觉语言模型}（vision-language model, VLM）是完成这种连接的一类模型，它把视觉输入与语言联系起来，可以输出可比较表示，也可以生成语言。本书采用这一较宽的用法；连接音频、深度或惯性信号的模型还需相应编码和训练依据，不能仅凭“多模态”名称确定支持范围。

连接发生的位置有三种安排，它们对“何时让两种模态相遇”给出不同回答。第一种安排是两侧独立编码，再把表示映射到共同空间计算相似度。图像向量不依赖当前文字查询，适合预先建立图库。第二种安排让文字计算读取视觉特征：语言隐状态提供查询，也就是当前希望取用信息的向量；视觉特征提供用于比较的键和承载内容的值，再按匹配权重读取。此时视觉底层特征可以缓存，但读取结果随问题和已生成文字改变。第三种安排把多源表示放入共享主干，通过位置、模态标记和注意力范围组织交互。

这三种安排是\ref{chap:vision-multimodal-alignment}章对齐与融合方法在大模型中的具体落点。第一种安排对应\ref{sec:vision-alignment-contrastive}节基于对比学习的对齐：该节说明了为什么用配对样本拉近、非配对样本推开就能形成共同空间，以及批内负例和温度起什么作用；本章的CLIP与ImageBind是这一原理的两个实现。第二种安排对应\ref{sec:vision-alignment-attention}节的注意力软对齐，以及\ref{sec:vision-alignment-intermediate-fusion}节以交叉注意力实现的中间融合：前者把注意力权重解释为两组单元之间的软对应，后者说明在网络中间层让一种模态调制另一种模态的计算。第三种安排则接近\ref{sec:vision-alignment-early-fusion}节的早期融合，只是融合的对象不再是手工特征的拼接，而是进入同一主干的词元序列。那一章讲各方法的原理和适用条件，本章不再重复，只关心这些原理在具体模型中取怎样的形状、怎样训练。

第二种安排的核心运算是\term{交叉注意力}（cross-attention），即一组表示提出查询、去另一组表示中加权读取内容；其计算在卷五已经讲过，作为对齐和融合手段的原理见\ref{sec:vision-alignment-attention}节，也可参考\BookChapterRef{chap:transformer}{卷五“Transformer”章}。这里把它写成一次可检查的读取，以明确哪里依赖当前问题。设语言隐状态为$\bm H\in\mathbb R^{m\times d_{\mathrm l}}$，视觉特征为$\bm F\in\mathbb R^{n_{\mathrm v}\times d_{\mathrm v}}$，分别投影得到$\bm Q\in\mathbb R^{m\times d_{\mathrm a}}$、$\bm K\in\mathbb R^{n_{\mathrm v}\times d_{\mathrm a}}$与$\bm V_{\mathrm a}\in\mathbb R^{n_{\mathrm v}\times d_{\mathrm o}}$。每个文字位置按与视觉位置的相关程度加权读取视觉值：
\begin{equation}
  \bm R=\operatorname{softmax}_{\mathrm{row}}
     \left(\frac{\bm Q\bm K^\top}{\sqrt{d_{\mathrm a}}}+\bm M\right)
     \bm V_{\mathrm a}.
  \label{eq:vision-models-cross-read}
\end{equation}
这里$\bm M\in\mathbb R^{m\times n_{\mathrm v}}$在允许读取的位置取0，禁止读取的位置取$-\infty$，每行至少要有一个可读位置。行softmax使一个文字位置的权重和为1，$\bm R\in\mathbb R^{m\times d_{\mathrm o}}$的每一行是相应视觉值的加权和。问题从“车是什么颜色”变成“车窗是否打开”时，$\bm H$及查询改变，即使$\bm F$未变，读取权重也可改变。注意力数值只表示内部加权，不自动构成对象定位的标注。

三种安排描述计算发生在哪里，不是互斥的完整模型类别。一个系统可以先用独立表示召回候选，再让联合模型逐个检查。还需区分三个常被混用的说法：表示对齐指训练使不同来源的内容在比较或交互中形成有用联系，参数共享指不同输入实际使用相同的权重，连接模块则负责把形状或信息路径接起来。两套编码器即使输出同样维数，也没有因此完成语义对齐。

三种安排的代价也不同。以图库查询为例，独立编码可以预先保存$M$幅图的向量，新请求只需编码一次文字并做$M$次内积。若每幅图都要与当前问题共同通过一个大型主干，就要执行$M$次联合计算。独立编码便宜，却只能比较两个整体向量，难以判断“左车红、右车蓝”这类需要逐个对象对照的关系；联合计算能逐位置读取，却难以在预算内处理全部候选。下一节的CLIP和ImageBind采用第一种安排，语言响应模型多采用第二、三种安排。

\subsection{输出表示的解码}
\label{sec:vision-models-output}

连接后的表示还不是用户要的结果。输出分数、输出文字和输出媒体需要不同的预测对象：匹配接口计算一个向量内积；语言接口在每一步给出词表上的分布，再逐词选出词元；图像或音频接口则预测像素、波形或能够被解码的中间表示。中间表示必须与解码器的训练约定相容，任意同维向量都不能送入同一个解码器。

表\ref{tab:vision-model-interfaces}固定同一街景输入，比较不同请求所需的结果。若给一个图文匹配模型附加图像生成器，匹配模型可以提供条件，但生成能力属于完整连接后的系统。反过来，视觉问答虽然输出文字，评价的仍是回答是否依据画面作出了正确判断。后面三节依次讨论第一行（理解模型）、后三行（生成模型）以及把两类输出放进同一系统的做法。

\begin{table}[htbp]
  \centering
  \small
  \begin{tabularx}{\linewidth}{@{}p{22mm}XX@{}}
    \toprule
    请求 & 模型直接预测的对象 & 恢复最终结果还需什么\\
    \midrule
    选择车辆类别 & 候选描述的匹配分数 & 给定候选集与选择规则\\
    回答停车位置 & 下一文字词元的分布 & 逐词解码与来源核验\\
    生成雨天变体 & 图像潜表示的更新量 & 生成迭代与图像解码器\\
    说出判断 & 语音编码序列 & 声学解码与波形合成\\
    \bottomrule
  \end{tabularx}
  \caption{输入相同而输出接口不同的教学比较。每行描述一种完整路径，并不意味着同一个检查点同时具有四种能力。}
  \label{tab:vision-model-interfaces}
\end{table}
\FloatBarrier

\section{多模态理解模型}
\label{sec:vision-models-understanding}

本节的问题是：给定一种模态的输入，怎样在另一种模态的候选中找出与之对应的内容。例如用“红色汽车”这段文字在图库中找照片，或用一段发动机录音找画面。输入是两种模态的数据，输出是它们之间的相似度，进而得到排序或类别。难点在于两侧编码器的坐标原本没有对应关系，而人工标注“哪段文字对应哪张图”的成本很高；能否利用网络上大量自然出现的配对，就成为关键。

两个代表模型都采用上一节的独立编码安排，用对比目标把配对的两侧拉近、把非配对的推开，差别在于用哪些配对来训练、哪些编码器参与更新。CLIP直接用图文对，从头同时训练图像和文字两个编码器，得到一个图文共同空间。ImageBind要连接更多模态，若为每两种模态都收集配对，成本随模态数平方增长；它改为只收集“图像—其他模态”的配对，并冻结CLIP式的图像和文字编码器，让新模态向已经形成的图文空间靠拢。前者回答“怎样从零建立两种模态的共同空间”，后者回答“怎样以较少配对把更多模态接进已有空间”。

\subsection{CLIP}
\label{sec:vision-models-contrastive}

在CLIP之前，图像分类器在固定类别集上用人工标注训练，输出层的每一维对应一个类别。要识别训练时没有的新类别，就必须重新收集标注并训练输出层；类别名本身携带的语义，例如“汽车”和“卡车”很接近，也没有被利用。\term{对比语言图像预训练}（contrastive language-image pre-training, CLIP）改用网络上自然出现的图片及其说明文字作为监督，把“预测固定类别编号”换成“判断哪段文字与这幅图配对”。原论文收集了约4亿个图文对\scite{radford2021}。

主要想法是分别编码两种输入，再要求一个批次中真实配对的相似度高于其他所有组合。这样做的好处是：文字端可以接受任意描述，类别不再是固定编号，新类别只要能写成文字就能比较；而真实配对来自网络，规模远大于人工标注。代价是监督只到整图与整句的对应，不告诉模型哪个词对应哪块区域。

模型由图像编码器、文字编码器和两个线性投影组成，两侧网络不必相同。原论文训练了五个ResNet和三个ViT图像编码器，每个都配一个Transformer文字编码器，表\ref{tab:vision-models-clip-config}列出它们的配置。

ResNet一侧在\ref{sec:vision-ivm-resnet}节的瓶颈残差结构上做了三处修改：采用ResNet-D的改进、加入抗混叠的rect-2模糊池化，并把末端的全局平均池化换成\term{注意力池化}（attention pooling）。注意力池化是一层多头查询—键—值注意力，查询由全局平均池化后的整图向量形成，键和值来自最后一个阶段的空间特征；它的输出直接作为整图表示，相当于让整图向量按内容重新加权各位置，而不是平等平均。表中的“各阶段块数”如RN50的$(3,4,6,3)$，表示四个阶段分别重复多少个瓶颈块；RN50x4、RN50x16、RN50x64按EfficientNet式的规则同时加宽、加深并提高输入分辨率，计算量约为RN50的4、16、64倍。

ViT一侧沿用\ref{sec:vision-ivm-vit}节的结构，只在块嵌入与位置嵌入相加之后、进入Transformer之前多加一层层归一化。以ViT-L/14为例，$224\times224$的图像切成$16\times16=256$个$14\times14$像素块，加一个分类位置共257个单元，每个单元宽1024；24个Transformer块逐层更新这$257\times1024$的数组，每块含16个注意力头。单个块按先归一化后残差的顺序计算：
\begin{equation}
  \bm Z'=\bm Z+\operatorname{MSA}\bigl(\operatorname{LN}(\bm Z)\bigr),
  \qquad
  \bm Z''=\bm Z'+\operatorname{MLP}\bigl(\operatorname{LN}(\bm Z')\bigr),
  \label{eq:vision-models-transformer-block}
\end{equation}
其中$\operatorname{LN}$是逐位置的层归一化，$\operatorname{MSA}$是多头自注意力，$\operatorname{MLP}$是逐位置的两层前馈网络，形状始终保持$257\times1024$。论文没有单独说明整图向量取自哪个位置；按ViT的做法，由最后一层分类位置的输出给出整图向量$f_{\mathrm v}(\bm X)\in\mathbb R^{1024}$。

文字编码器也是由式\eqref{eq:vision-models-transformer-block}的块堆叠而成的Transformer，基础规模为12层、宽512、8个注意力头，约6300万参数；图像编码器变大时，文字编码器只按比例加宽，层数不变。文字先转为小写，用字节对编码切成词元，超参数表中的词表大小为49408；序列首尾加起始和结束标记，最长76个词元。自注意力使用因果掩码，每个位置只读取它之前的词元，这保留了日后用预训练语言模型初始化的可能。结束标记处最高层的隐状态经层归一化，作为整句向量$f_{\mathrm t}(q)$。

\begin{table*}[tbp]
  \centering
  \small
  \begin{tabularx}{\linewidth}{@{}lcXcccc@{}}
    \toprule
    图像编码器  & 输入分辨率  & 结构  & 视觉宽度  & 视觉头数  & 文字层数/宽度/头数  & 嵌入维数$d$\\
    \midrule
    RN50  & 224  & 各阶段块数$(3,4,6,3)$  & 2048  & —  & 12/512/8  & 1024\\
    RN101  & 224  & 各阶段块数$(3,4,23,3)$  & 2048  & —  & 12/512/8  & 512\\
    RN50x4  & 288  & 各阶段块数$(4,6,10,6)$  & 2560  & —  & 12/640/10  & 640\\
    RN50x16  & 384  & 各阶段块数$(6,8,18,8)$  & 3072  & —  & 12/768/12  & 768\\
    RN50x64  & 448  & 各阶段块数$(3,15,36,10)$  & 4096  & —  & 12/1024/16  & 1024\\
    \midrule
    ViT-B/32  & 224  & 块32，12层，50个单元  & 768  & 12  & 12/512/8  & 512\\
    ViT-B/16  & 224  & 块16，12层，197个单元  & 768  & 12  & 12/512/8  & 512\\
    ViT-L/14  & 224  & 块14，24层，257个单元  & 1024  & 16  & 12/768/12  & 768\\
    ViT-L/14@336  & 336  & 块14，24层，577个单元  & 1024  & 16  & 12/768/12  & 768\\
    \bottomrule
  \end{tabularx}
  \caption{CLIP各图像编码器与配套文字编码器的配置，数值取自原论文附录的超参数表。ViT的单元数由$(\text{分辨率}/\text{块大小})^2+1$算得，含一个分类位置；ResNet的“视觉宽度”为末阶段通道数，注意力池化的头数论文未单列。ViT-L/14@336是在ViT-L/14基础上以336像素再训练一轮的版本。各模型的参数量与MLP倍率论文未给出。}
  \label{tab:vision-models-clip-config}
\end{table*}

两侧编码器之后各接一个线性投影，论文没有使用非线性投影头。以ViT-L/14为例，图像侧把1024维映射到768维，文字侧把768维映射到768维。设一批有$n$对图文$(\bm{X}_i,q_i)$，两个编码器分别输出$f_{\mathrm v}(\bm{X}_i)\in\mathbb R^{d_{\mathrm v}}$与$f_{\mathrm t}(q_i)\in\mathbb R^{d_{\mathrm t}}$，投影矩阵$\bm{W}_{\mathrm v}\in\mathbb R^{d\times d_{\mathrm v}}$、$\bm{W}_{\mathrm t}\in\mathbb R^{d\times d_{\mathrm t}}$将它们映射到$\mathbb{R}^d$。对于非零投影结果，归一化表示为
\begin{equation}
  \bm{v}_i=\frac{\bm{W}_{\mathrm v}f_{\mathrm v}(\bm{X}_i)}
                    {\lVert\bm{W}_{\mathrm v}f_{\mathrm v}(\bm{X}_i)\rVert_2},
  \qquad
  \bm{u}_i=\frac{\bm{W}_{\mathrm t}f_{\mathrm t}(q_i)}
                    {\lVert\bm{W}_{\mathrm t}f_{\mathrm t}(q_i)\rVert_2}\text{。}
  \label{eq:vision-clip-embedding}
\end{equation}
内积$s_{i,j}=\bm{v}_i^\top\bm{u}_j$就是余弦相似度，$n$对图文形成$n\times n$的相似度矩阵$\bm S$。归一化使向量长度不能单独抬高分数，模型必须改变方向来改善配对。

学习目标把相似度矩阵变成两个分类问题。温度$\tau>0$把相似度转换为$\ell_{i,j}=s_{i,j}/\tau$，这些未归一化分类分数称为logit；温度越小，同样的相似度差会形成越大的logit差。固定一幅图，在本批次$n$段文字中识别其配对文本，使用行方向交叉熵，即为正确配对的负对数选择概率计分；固定一段文字，反过来在图像中识别配对，则使用列方向交叉熵：
\begin{equation}
  \begin{aligned}
    L_{\mathrm{v\to t}}
      &=-\frac1n\sum_{i=1}^n
        \log\frac{\exp(\ell_{i,i})}{\sum_{j=1}^n\exp(\ell_{i,j})},\\
    L_{\mathrm{t\to v}}
      &=-\frac1n\sum_{i=1}^n
        \log\frac{\exp(\ell_{i,i})}{\sum_{j=1}^n\exp(\ell_{j,i})},\\
    L_{\mathrm{CLIP}}&=\tfrac12
        \bigl(L_{\mathrm{v\to t}}+L_{\mathrm{t\to v}}\bigr)\text{。}
  \end{aligned}
  \label{eq:vision-clip-loss}
\end{equation}
对角线表示数据记录中的正确配对，其他单元在该目标中充当负例。两项损失约束相同的相似度矩阵，但分母分别沿行和列求和。若批次中另有一幅也符合该描述的图，单一配对标签会把它当负例，这就是配对监督和真实语义关系之间的差别。对比目标为什么能把配对拉近、把非配对推开，批内负例数与温度怎样影响学习信号，已在\ref{sec:vision-alignment-contrastive}节讲过，更一般的对比学习可参考\BookChapterRef{chap:data-derived-knowledge-acquisition}{卷八“基于数据本身的知识获取”章}；这里只关注CLIP采用的对称形式，以及图文两侧怎样共同构成配对判定。

交叉熵为何会改变相对得分，可以从一行的目标直接推导。令$p_{i,j}=\exp(\ell_{i,j})/\sum_k\exp(\ell_{i,k})$，第$i$行的误差为$-\ell_{i,i}+\log\sum_k\exp(\ell_{i,k})$。第一项只在$j=i$时导数为$-1$，第二项的导数为$p_{i,j}$，两者相加得到$p_{i,j}-\mathbb I[j=i]$。若再令$\widetilde p_{i,j}$为第$j$列沿图像方向归一化的概率，完整双向平均的梯度为
\begin{equation}
  \frac{\partial L_{\mathrm{CLIP}}}{\partial\ell_{i,j}}
   =\frac{p_{i,j}+\widetilde p_{i,j}
                 -2\mathbb I[i=j]}{2n}.
  \label{eq:vision-clip-logit-gradient}
\end{equation}
正配对概率还小于1时，对角线的导数为负；负配对的导数为正。沿负梯度改变参数，便倾向于增大前者、减小后者。logit不是独立可调参数，实际更新还须通过内积、归一化和两个编码器传回权重。对投影结果$\bm a\ne\bm0$与$\bm v=\bm a/\lVert\bm a\rVert_2$，归一化的导数矩阵为$(\bm I-\bm v\bm v^\top)/\lVert\bm a\rVert_2$，它去掉沿$\bm v$方向、只改变长度的梯度分量，保留能改变方向的分量。图像编码器收到同一行所有文字比较的梯度，文字编码器收到同一列所有图像比较的梯度，所以一次真实配对同时约束了两侧。

取一个可复算的二维教学批次：两幅图的表示分别为$(1,0)^\top$、$(0,1)^\top$，两段配对文字“红色汽车”“红色自行车”的表示分别为$(0.8,0.6)^\top$、$(0.6,0.8)^\top$，四个向量均为单位向量。取$\tau=0.2$，得到
\[
  \bm{S}=\begin{bmatrix}0.8&0.6\\0.6&0.8\end{bmatrix},
  \qquad
  \bm{\ell}=\begin{bmatrix}4&3\\3&4\end{bmatrix}\text{。}
\]
每一行的正确配对概率为$e^4/(e^4+e^3)\approx0.7311$，交叉熵为$\log(1+e^{-1})\approx0.3133$；列方向在这个对称例子中相同。完整双向平均中$n=2$，式\eqref{eq:vision-clip-logit-gradient}给出第一行导数$(-0.1345,0.1345)$。若暂把logit当独立变量、取教学步长0.1，$(4,3)$更新为约$(4.0134,2.9866)$，正确选择概率增至约0.7363。这不是真实网络的一次训练结果，它只复算目标希望把两个分数差推向哪个方向。图\ref{fig:vision-clip-matrix}标出了两个归一化方向。

\begin{figure}[htbp]
  \centering
  % 教学数据由四个单位向量的内积构造，温度为0.2。
\begin{tikzpicture}[x=1cm,y=1cm,font=\figurefont\small,
  arr/.style={-{Stealth[length=2mm]},BookInk,line width=0.8pt}]
  \node[anchor=west] at (0,4.5) {缩放后分数$\ell_{i,j}=s_{i,j}/\tau$};
  \node[align=center] at (3.8,3.55) {$q_1$\\红色汽车};
  \node[align=center] at (6.0,3.55) {$q_2$\\红色自行车};
  \node[anchor=east] at (2.6,2.5) {图像1};
  \node[anchor=east] at (2.6,1.2) {图像2};
  \foreach \x/\y/\val in {3.8/2.5/4,6.0/2.5/3,3.8/1.2/3,6.0/1.2/4}{
    \draw[BookRule,line width=0.5pt,fill=BookPaper]
      (\x-1,\y-0.55) rectangle (\x+1,\y+0.55);
    \node at (\x,\y) {$\val$};
  }
  \draw[BookTeal,line width=1.2pt] (2.8,1.95) rectangle (4.8,3.05);
  \draw[BookTeal,line width=1.2pt] (5.0,0.65) rectangle (7.0,1.75);
  \draw[arr] (7.2,2.5) -- (8.15,2.5);
  \node[anchor=west,align=left] at (8.2,2.5) {行softmax\\$0.7311$\\$0.2689$};
  \draw[arr] (3.8,0.4) -- (3.8,-0.2);
  \node[align=center,anchor=north] at (3.8,-0.3)
    {列softmax\\$(0.7311,\;0.2689)^\top$};
  \node[anchor=west] at (0,-1.35) {对角线：数据配对\qquad 非对角线：批次负例};
\end{tikzpicture}

  \caption{CLIP双向目标的两对图文教学示例。行固定图像、列固定文字，粗框标出训练配对。行、列softmax分别形成两个分类问题，图中的0.7311是批次内选择概率，不是“图中确有该物体”的校准概率。}
  \label{fig:vision-clip-matrix}
\end{figure}

训练时，先把图片按模型约定缩放后随机裁出正方形区域（这是唯一的数据增强），把对应文本按固定分词器编码，再抽取一批记录构造整个$\bm S$。原论文使用32768的批次大小，使每个正配对在一次更新中同时与三万多个负例比较。CLIP从随机初始化的编码器开始，用带解耦权重衰减的Adam优化器（衰减系数0.2，学习率先线性预热2000步再按余弦日程下降）共同调整两个编码器、投影和温度，所有模型都训练32轮，按这一日程停止，而不是等到全部配对概率变成1；最大的RN50x64在592块V100上训练了18天，最大的ViT在256块V100上训练了12天。温度通过学习对数尺度保持为正，初值相当于$\tau=0.07$，并截断使$1/\tau$不超过100，避免logit被放得过大导致训练不稳定；它改变softmax的尖锐程度，不改变未缩放余弦分数的排序。训练后保留图像处理方式、分词器、两侧权重和投影，丢弃当前批次的负例矩阵。下降的训练目标没有给出新数据上的正确率保证，后者仍需另行评价。

使用时，可以把类别$c$写成描述模板$q_c$，缓存其表示$\bm{u}_c$，再对新图像选取$\argmax_c\bm{v}^\top\bm{u}_c$。这种零样本分类不使用当前类别任务的标注样本去拟合分类器，但仍依赖大规模图文预训练。写“这是一张红色汽车的照片”和只写“汽车”，提供给文字编码器的上下文不同，得到的方向也可能不同；若候选只有“汽车”“自行车”两个，模型一定会从中选一个，除非另加拒绝规则。比较结果时应记录候选集合、模板、图像处理和模型版本。整图相似度没有直接给出车辆的框、实例掩码或自由答案，这些输出需要其他接口。

\subsection{ImageBind}
\label{sec:vision-models-imagebind}

图文配对之外，视频帧和同时录到的声音、彩色图与对应深度、佩戴设备的画面与惯性信号，也提供了跨来源联系。若沿CLIP的做法为每两种模态都建立共同空间，就需要每两种模态之间的配对数据：六种模态有15种组合，而“热成像—文字”“惯性信号—音频”这类配对几乎不存在。ImageBind的主要想法是用图像作为配对中心：每种新模态只需与图像配对，就把它对齐到图像所在的共同空间；由于图像已经通过CLIP式训练与文字对齐，新模态与文字之间也可能形成联系，尽管二者从未直接配对\scite{vision-imagebind2023}。

模型为六种模态各设一个Transformer编码器，表\ref{tab:vision-models-imagebind-config}按组件列出论文公开的配置。图像与视频共用视觉编码器，它是OpenCLIP的ViT-H，约6.3亿参数，块大小$16\times16$；视频输入取短时间范围内的两帧，把块投影沿时间“膨胀”为$2\times16\times16$的时空块，于是单幅图像和两帧视频可以进入同一个编码器。文字编码器沿用CLIP的文字编码器设计，取自同一OpenCLIP模型，约3.02亿参数。这两个编码器在ImageBind训练中都冻结。

其余四种模态的编码器都从头训练，结构是式\eqref{eq:vision-models-transformer-block}的Transformer块堆叠。音频以16kHz采样，取约2秒的片段，按上一节的方法用25毫秒汉明窗、10毫秒步长和128个梅尔频带转为频谱，再像图像一样用边长16、步长10的块切分；块之间有重叠，这与ViT的不重叠切块不同。按这些参数估算，2秒约有200个时间段，频率方向可切$\lfloor(128-16)/10\rfloor+1=12$个块位置，时间方向约19个，合计约230个音频单元。深度先换成单通道视差图，热成像也作为单通道图像，二者都按图像方式切块。惯性测量单元的信号是加速度计与陀螺仪各三轴的六通道序列，取5秒片段共约2000个时间步，用核长为8的一维卷积投影成序列，再送入Transformer。各编码器末端各有一个线性投影，把输出映射到同一维数$d$并归一化后参与对比目标；论文比较过多层感知机投影头，线性头更好，但$d$的具体数值论文未给出。

\begin{table}[htbp]
  \centering
  \small
  \begin{tabularx}{\linewidth}{@{}p{14mm}XXp{12mm}@{}}
    \toprule
    组件 & 输入与输出 & 结构与规模 & 训练状态\\
    \midrule
    图像/视频 & 图像或两帧片段，$16\times16$块（视频$2\times16\times16$） & OpenCLIP ViT-H，约6.3亿参数 & 冻结\\
    文字 & 词元序列 & OpenCLIP文字编码器，约3.02亿参数 & 冻结\\
    音频 & 约2秒、128频带梅尔频谱，块16步长10 & 12层，宽768，12头 & 训练\\
    深度 & 单通道视差图 & 12层，宽384，8头 & 训练\\
    热成像 & 单通道图像 & 12层，宽768，12头 & 训练\\
    惯性信号 & $6\times2000$序列，核长8的一维卷积投影 & 6层，宽512，8头 & 训练\\
    投影头 & 各编码器输出$\to d$维 & 每模态一个线性层，$d$论文未给出 & 训练\\
    \bottomrule
  \end{tabularx}
  \caption{ImageBind各模态编码器的配置。四种新模态的层数、宽度和头数取自论文附录的预训练超参数表，该表按配对数据集列出（音频对应AudioSet，深度对应SUN RGB-D，热成像对应LLVIP，惯性信号对应Ego4D）；冻结编码器的层数与宽度论文未单列。}
  \label{tab:vision-models-imagebind-config}
\end{table}

学习方法沿用CLIP的双向目标，只是配对换成“图像—新模态”。设模态$m$的数据集为$\mathcal D_m=\{(\bm X_i,a_i^{m})\}$，例如视频帧与同时录到的声音，其编码与投影参数为$\bm\theta_m$。每次在一个配对数据集中抽取批次，计算图像向量$\bm v_i$与该模态向量$\bm u^{m}_j$的相似度矩阵，按式\eqref{eq:vision-clip-loss}形成双向目标$L_m$。原论文冻结图像和文字编码器，误差只经归一化和投影反传到$\bm\theta_m$，于是新模态必须去适配图像空间，而不是让所有编码器同时任意移动；文字端也不因音频批次改变。各模态分别训练，不要求一条记录同时拥有全部模态。与CLIP学习温度不同，ImageBind的最佳设置对每个配对数据集固定温度：音频—视频取0.05，图像—深度取0.2，图像—热成像取0.1，视频—惯性信号取0.2；论文的消融显示，从0.07起学习温度不如固定值。各配对用AdamW训练，音频批次为2048，其余为512，惯性信号训练8轮，其余64轮。训练后保留各自的预处理、编码器与投影；推理不再需要配对标签和批次内负例。

使用时，用声音检索画面只需编码声音，再按它与图库向量的内积排序。继续使用街景例子，假设发动机录音形成单位向量$(0.9,\sqrt{0.19})^\top$，图库中两幅候选的图像表示为$(1,0)^\top$和$(0,1)^\top$，两个内积约为0.9和0.4359，排序先返回前者。若两幅候选都含车辆但只有后者是这段录音的真实来源，几何排序仍可能选错；训练配对学到的是所覆盖数据中的联系，并非每次录音的声源认证。用声音选择文字类别时，可以比较声音向量与候选文字向量，即使音频编码器没有用音文对训练。

这种未直接配对模态之间的联系为什么可能出现，可以用一个几何事实说明。若某段声音和某段文字都与同一个图像向量很接近，即$\lVert\bm{u}^{\mathrm a}-\bm v\rVert_2\leq\epsilon_{\mathrm a}$、$\lVert\bm{u}^{\mathrm t}-\bm v\rVert_2\leq\epsilon_{\mathrm t}$，三角不等式给出
\[
  \lVert\bm{u}^{\mathrm a}-\bm{u}^{\mathrm t}\rVert_2
  \leq\epsilon_{\mathrm a}+\epsilon_{\mathrm t}\text{。}
\]
但实际训练不能让每个语义一致的样本都围绕同一个准确图像锚点，也不能消除背景、视角和声源歧义，因此这只是有条件的几何解释。论文在所测数据中观察到了这种联系；共同空间本身不是保证联系正确的定理。画面中既有车又有树而录音只有风声时，整段音画配对也没有显式指出声音来自哪个区域。ImageBind的音频到图像生成展示还外接了图像生成系统，用声音表示替代相应条件；ImageBind本体仍只输出表示。

\subsection{理解结果的形成与接口选择}
\label{sec:vision-models-understanding-interface}

CLIP与ImageBind都把理解结果表示为共同空间中的向量，判断靠内积排序完成。表\ref{tab:vision-models-understanding-compare}在同一检索任务下比较二者。ImageBind改变了CLIP的两项设计：配对从图文扩展为以图像为中心的多种模态，图像和文字编码器从共同训练改为冻结。换来的是不需要两两配对就能接入新模态；付出的代价是新模态只能适应已经固定的图文空间，若某种声音在图文数据中缺少相应概念，冻结空间里就没有合适的位置，而且未直接配对模态之间的检索质量依赖图像这一中介。

\begin{table}[htbp]
  \centering
  \small
  \begin{tabularx}{\linewidth}{@{}p{20mm}XX@{}}
    \toprule
    比较项 & CLIP & ImageBind\\
    \midrule
    训练配对 & 图文对 & 图像与音频、深度、热成像、惯性信号等各自配对\\
    更新的参数 & 图像、文字编码器与投影共同训练 & 图文编码器冻结，只训练新模态编码器与投影\\
    覆盖的模态 & 图像、文字 & 六种模态共用一个空间\\
    未直接配对的比较 & 不涉及 & 经图像中介形成，属于经验观察\\
    输出 & 整图与整句向量、相似度 & 各模态整体向量、相似度\\
    \bottomrule
  \end{tabularx}
  \caption{两个多模态理解模型的比较。两者都只输出可比较表示，不输出区域或文字。}
  \label{tab:vision-models-understanding-compare}
\end{table}

相似度接口适合在已定义候选中选择，却不能回答所有理解请求。街景查询“找相似画面”可以直接返回向量排序；“描述发生了什么”没有固定候选句表，需要根据视觉内容预测一串文字；“指出右侧车辆”则要求把表达对应到区域，仅有一句“右边有一辆车”尚未给出可检查的位置。三者都涉及理解，但需要的结果不同：相似度适合选择，语言生成适合组织可变长度的描述和答案，定位接口则要输出框、点或掩码。语言模型也可以按约定输出坐标，但仍需训练依据、坐标转换和定位评价；语言指代与区域对应见\ref{sec:vision-understanding-grounding}节。

因此，接口选择先看输出形式。候选可枚举且只需排序时，选择双编码表示，它可以预先缓存图库向量，成本最低；需要自由描述或回答时，选择下一节的语言响应模型；需要位置时，选择带定位头的模型。同一表示也可服务生成：图像语义向量可以作为生成器的条件，图文相似度的梯度可以调整采样方向，生成结束后的相似度还可以筛选候选。这些用途分别改变生成输入、生成过程和选择规则，不能混为“CLIP生成图像”。下一节从预测对象和训练模块出发，解释这些不同连接。

\section{多模态生成模型}
\label{sec:vision-models-generation}

理解模型只给出候选之间的相似度，许多请求却要求模型产生新的内容：用一句话回答问题，画出雨天街景，或生成一段带声音的视频。这些输出的共同点是长度或尺寸可变、事先没有候选列表，模型必须逐步构造结果。难点随输出类型而不同。语言响应主要需要把已观察到的内容表达出来，困难在于让语言计算真正读取视觉证据；图像需要决定每处颜色、纹理与空间关系，困难在于高维像素的生成成本与条件控制；视频和音频还要在时间上保持连贯，并让画面和声音互相对应。

因此本节按输出类型分为三组。语言响应生成以已有的大型语言模型为输出端，各模型的差别在于视觉信息从哪里、以多少单元进入语言计算；图像生成以扩散模型为基础，各模型的差别在于在哪个空间生成、条件怎样进入去噪网络；视频与音频生成以流匹配训练的生成主干为基础，差别在于只生成画面还是让画面和声音共同生成。三组可以在同一系统中结合，第\ref{sec:vision-models-unified}节再讨论这种结合。

\subsection{语言响应生成}
\label{sec:vision-models-language-generation}

对街景问题“右侧停着什么车”，需要从同一视觉证据形成长度可变的答案，而不是从两个类别描述中选一个。输入是视觉表示$\bm{F}$和请求$q$，输出是回答的词元$y_1,\ldots,y_m$。条件语言生成按此前已经产生的前缀逐词展开：
\begin{equation}
  p_{\bm\theta}(y_1,\ldots,y_m\mid\bm F,q)
  =\prod_{j=1}^m p_{\bm\theta}(y_j\mid y_{<j},\bm F,q)\text{。}
  \label{eq:vision-vlm-language}
\end{equation}
这种自回归分解和因果掩码在卷五Transformer章已经讲过；这里的新问题是，式中的条件依赖只有通过实际连接才能成立。从头训练一个同时懂图像和语言的大模型代价极高，而已有大型语言模型已经掌握了语言知识，已有视觉编码器也已经学会提取视觉特征。本小节各模型的共同出发点，是复用这两个已经训练好的部件，只设计和训练两者之间的连接。

训练记录包含图片、请求和参考回答，必要时再含先前轮次及其角色。令$\mathcal T_i$为第$i$条记录中需要监督的回答位置，用户问题和作为条件的图片不属于这组目标位置。参考序列的概率是逐词条件概率的乘积，取负对数后得到可相加的误差；对$n$条记录形成经验目标
\begin{equation}
  \widehat L_{\mathrm{text}}
   =-\frac1n\sum_{i=1}^n\sum_{j\in\mathcal T_i}
      \log p_{\bm\theta}(y_{i,j}\mid
          y_{i,<j},\bm F_i,q_i).
  \label{eq:vision-vlm-text-objective}
\end{equation}
若参考回答中的“汽车”在当前位置预测概率为0.2，这个词元贡献$-\log0.2\approx1.6094$；若概率增到0.5，贡献降为约0.6931。负对数概率促使模型提高标注词元的概率，不直接判断标注句子的事实真假。使用时，语言头把当前位置隐向量映射为$V$个logit，softmax后得到词元分布；可以取最大概率词元（贪心解码），按分布抽样，或保留若干前缀按累计对数概率选择（束搜索）。选出的词元接回上下文，再计算下一位置，直到生成结束标记或达到长度上限。

在这个共同框架下，各模型对“视觉信息怎样到达预测$y_j$所用的语言隐状态”给出不同回答。BLIP-2先用少量可学习查询把视觉特征压缩成固定个数的向量，再放到语言输入前面；Flamingo不改变语言输入，而在语言模型的层间插入读取视觉特征的交叉注意力，便于处理图文交错的长上下文；LLaVA不压缩，直接把每个视觉网格位置线性投影到语言嵌入空间。这三种做法依次对应“前缀压缩”“层间读取”“直接投影”。Qwen3-VL沿用直接投影，但进一步把视觉编码器的多个中间层接入语言层，并为视频加入时间组织；Qwen2.5-Omni把输入扩展到音频，输出扩展到语音；CounterCredit则不再假设一次编码就够，让模型在回答过程中决定是否再次取回图像区域。前五个模型改变的是连接结构，最后一个改变的是观察何时发生。

\subsubsection{BLIP-2}
\label{sec:vision-models-blip2}

图像编码器为一幅图输出数百个特征，已有语言模型却没有学过怎样读取它们。若把全部特征直接送入语言模型，序列长度和计算量都很大；若只取一个整图向量，又会丢掉大部分细节。更根本的问题是，冻结的语言模型和冻结的图像编码器之间存在表示差距，一个简单投影很难让语言模型理解图像特征。BLIP-2的主要想法是在二者之间训练一个轻量的\term{查询Transformer}（querying transformer, Q-Former），即用一组可学习查询从图像特征中抽取信息的网络，并分两个阶段训练：先在不接语言模型的情况下让查询学会提取与文字相关的视觉信息，再把查询输出接入语言模型\scite{vision-blip22023}。

查询在这里是模型参数，不是用户的问题。设冻结图像编码器的输出为$\bm F\in\mathbb R^{n_{\mathrm v}\times d_{\mathrm v}}$，论文试了两种冻结图像编码器：CLIP的ViT-L/14和EVA-CLIP的ViT-g/14，都去掉最后一层、取倒数第二层的输出。预训练输入为$224\times224$，块大小14，于是两者都输出$16\times16+1=257$个位置；ViT-L/14的特征为$257\times1024$，ViT-g/14的特征宽度论文未给出。查询矩阵为$\bm Q_0\in\mathbb R^{32\times768}$，即32个768维查询。Q-Former的每层中，查询先通过自注意力交换信息，再以式\eqref{eq:vision-models-cross-read}的交叉注意力读取$\bm F$，得到$32\times n_{\mathrm v}$的读取权重；输出$\bm Z\in\mathbb R^{32\times768}$。无论图像编码器输出多少特征，交给下一模块的查询输出数量保持32个。这种压缩减少后续语言计算所读的视觉单元，也要求查询选择并保留与任务相关的信息。以两个视觉值$(1,0)^\top$、$(0,1)^\top$及单个教学查询为例，若读取权重为$(0.25,0.75)$，输出贡献就是$(0.25,0.75)^\top$；这是信息汇聚，不是自动裁出右车。Q-Former还有一个与查询共享自注意力层的文字分支，可以读取文字输入；共享层从预训练BERT初始化，新增交叉注意力随机初始化。

Q-Former的骨架是BERT-base，即12层、宽768、12个注意力头的Transformer编码器\scite{devlin2019}。它内部有两个共用自注意力层的子模块：图像子模块处理32个查询，文字子模块处理文字词元。图像子模块的每一层先做自注意力，在允许时查询还能读取文字；每隔一层在自注意力之后插入一个交叉注意力，读取冻结图像特征$\bm F$；最后是逐位置前馈层。12层中因此有6层带交叉注意力。若把第$\ell$层的查询状态记为$\bm Z^{(\ell)}\in\mathbb R^{32\times768}$，带交叉注意力的一层按“自注意力$\to$读取$\bm F$的交叉注意力$\to$前馈”的顺序更新，每个子层都带残差连接与层归一化，形状始终是$32\times768$；不带交叉注意力的一层省去中间一步。交叉注意力的键和值由$\bm F$投影而来，所以无论图像特征是1024维还是更宽，查询一侧的宽度都保持768。论文给出Q-Former共约1.88亿参数，表\ref{tab:vision-models-blip2-config}按组件汇总这一数据流。

\begin{table}[htbp]
  \centering
  \small
  \begin{tabularx}{\linewidth}{@{}p{14mm}XXp{12mm}@{}}
    \toprule
    组件 & 输入与输出 & 结构与规模 & 训练状态\\
    \midrule
    图像编码器 & $224\times224$图像$\to257\times d_{\mathrm v}$（ViT-L/14为$257\times1024$） & CLIP ViT-L/14或EVA-CLIP ViT-g/14，取倒数第二层 & 冻结\\
    Q-Former & 32个查询与$257\times d_{\mathrm v}$特征$\to32\times768$ & BERT-base骨架，12层、宽768、12头，每隔一层插入交叉注意力；约1.88亿参数 & 训练\\
    投影 & $32\times768\to32\times d_{\mathrm l}$ & 一个全连接层 & 训练（第二阶段）\\
    语言模型 & 视觉前缀与文字$\to$回答词元 & OPT-2.7B/6.7B（仅解码器）或FlanT5-XL/XXL（编码器—解码器） & 冻结\\
    \bottomrule
  \end{tabularx}
  \caption{BLIP-2的组件与数据流。$d_{\mathrm v}$为图像特征宽度，$d_{\mathrm l}$为所选语言模型的嵌入宽度。预训练共用1.29亿幅图像；第一阶段250K步，批次2320（ViT-L）或1680（ViT-g），第二阶段80K步，批次1920（OPT）或1520（FlanT5），均用AdamW、峰值学习率$10^{-4}$与余弦衰减。}
  \label{tab:vision-models-blip2-config}
\end{table}

第一阶段连接冻结图像编码器和Q-Former，使用图文对同时学习三种目标，它们共用Q-Former参数，只用不同的注意力掩码区分。对比目标让图像查询和文字分别编码，二者在自注意力中不能互看，以免把待匹配文本直接泄漏给图像表示；设归一化查询输出为$\bm z_{i,k}$、文字表示为$\bm t_j$，图文分数取$\max_{1\leq k\leq 32}\bm z_{i,k}^{\top}\bm t_j$，再按式\eqref{eq:vision-clip-loss}形成批次目标。配对判定目标允许查询和文字双向交互，由各查询输出的二分类logit取平均，判断图文是否匹配；真实配对标签$b=1$、困难负例标签$b=0$，误差为$-b\log p_1-(1-b)\log p_0$。图像条件文字生成目标让查询之间互看但不能读取目标文字，每个文字词元可以读取全部查询和此前文字；由于文字不能直接读取图像特征，信息只能经查询传递，这迫使查询提取对生成描述有用的内容。三个误差相加，梯度传回共享层、查询和各目标的输出层，图像编码器权重不变。

第二阶段用$\bm W\in\mathbb R^{d_{\mathrm l}\times 768}$把$\bm Z$线性投影到语言模型的嵌入维数$d_{\mathrm l}$，得到$\bm P=\bm Z\bm W^\top\in\mathbb R^{32\times d_{\mathrm l}}$。这些向量放在文字输入前面，称为\term{视觉前缀}（visual prefix）：它们占据语言模型的输入位置，用连续向量提供图像条件，却不对应可直接打印的文字。冻结图像编码器和语言模型（OPT或FlanT5），用式\eqref{eq:vision-vlm-text-objective}训练Q-Former与投影层；误差经冻结语言层反传到$\bm P$，再更新投影与Q-Former。对编码器—解码器式语言模型，论文还使用前缀语言建模，把部分文字与视觉前缀交给编码器，预测剩余文字。原论文第一阶段训练250K步、第二阶段80K步，这些停止点是训练配置，不是证明目标已全局最小；图文描述部分由论文的描述生成与筛选流程获得，合成文字不应被当作无误标注。

在街景问答中，模型先计算图像特征，再取得32个查询输出，接上“右侧停着什么车”，由语言模型逐词生成答案。固定图片的视觉前缀可用于不同文字请求，但压缩时未保留的细节无法仅靠换一种问法恢复。这里的冻结说明针对两阶段预训练，不能推广到原论文所有下游微调设置。

\subsubsection{Flamingo}
\label{sec:vision-models-flamingo}

BLIP-2一次处理一幅图和一段文字。若输入是“示例图像—示例回答—待答图像—问题”这样的图文交错序列，视觉信息还需要与文字在序列中的位置对应：回答待答图像的问题时，应当读取待答图像，而不是示例图像。把所有图像的视觉前缀都拼进语言输入，序列会随图像数量快速变长，也难以表达“哪段文字对应哪幅图”。Flamingo的主要想法是不改变语言模型的输入序列，而在冻结语言层之间插入新的交叉注意力层，让每个文字位置只读取与它对应的图像；同时把每幅图像或视频压缩为固定数量的视觉单元，使读取成本不随原始特征数增长\scite{vision-flamingo2022}。

架构有两个新增模块。第一个是\term{感知重采样器}（Perceiver Resampler），它用64个可学习查询，通过注意力把冻结视觉编码器输出的数量可变的特征汇聚为64个输出；视频则对采样帧分别编码、加入时间嵌入后一起汇聚。这里的“重采样”是学习式信息汇聚，不是从原图均匀抽取64个像素；其注意力的键和值还包含查询隐表示本身。第二个是门控交叉注意力层：查询来自当前语言隐状态，键和值来自重采样后的视觉表示。若用$C(\bm H,\bm Z)$概括一个新增交叉注意力分支，残差更新的核心形式为
\begin{equation}
  \bm H'=\bm H+\tanh(\alpha)\,C(\bm H,\bm Z)\text{，}
  \label{eq:vision-flamingo-gate}
\end{equation}
其中$\alpha$是从0开始学习的标量门控。初始时新增分支对输出的贡献为0，语言模型保留原来的计算，训练不会一开始就破坏已有语言能力。完整的门控交叉注意力层（论文称gated xattn-dense）还有一个同样带门控的前馈分支，两步依次为
\begin{equation}
  \bm H'=\bm H+\tanh(\alpha_{\mathrm{x}})\,\operatorname{XAttn}(\bm H,\bm Z),
  \qquad
  \bm H''=\bm H'+\tanh(\alpha_{\mathrm{d}})\,\operatorname{FFW}(\bm H'),
  \label{eq:vision-flamingo-block}
\end{equation}
其后才是原语言模型冻结的自注意力与前馈层。式\eqref{eq:vision-flamingo-gate}只是其中第一步，$\alpha_{\mathrm x}$与$\alpha_{\mathrm d}$都从0开始学习。

感知重采样器自身是一个6层的Transformer。冻结的NFNet-F6视觉编码器（约4.35亿参数，Flamingo训练时输入$320\times320$）输出最后一个阶段的二维特征网格，展平成序列$\bm X_{\mathrm f}$；视频按每秒1帧取样，训练时取8帧，各帧独立编码后加上可学习的时间嵌入再拼接。64个可学习的潜查询$\bm X$在每层先读取再变换：
\[
  \bm X\leftarrow\bm X+\operatorname{Attn}\bigl(\bm X,[\bm X_{\mathrm f};\bm X]\bigr),
  \qquad
  \bm X\leftarrow\bm X+\operatorname{FFW}(\bm X),
\]
其中$\operatorname{Attn}(\cdot,\cdot)$的第一个参数提供查询，第二个参数沿序列方向拼接视觉特征与潜查询本身，提供键和值。6层之后输出64个视觉单元，这就是正文所说的$\bm Z$。由于查询数固定，交叉注意力每次读取的视觉单元数与帧数、分辨率无关。重采样器宽1536、16个头，三个规模共用同一结构，约1.94亿参数。

三个规模的差别在于冻结语言模型的大小和门控层插入的频率，见表\ref{tab:vision-models-flamingo-config}。Flamingo-3B在每个语言层前都插入门控层，9B每4层插入一层，80B每7层插入一层；语言模型越大，越稀疏地插入，以控制新增参数。新增层的前馈激活使用平方ReLU，冻结语言模型保持原来的GELU。

\begin{table}[htbp]
  \centering
  \small
  \begin{tabularx}{\linewidth}{@{}p{16mm}XXX@{}}
    \toprule
    组件 & Flamingo-3B & Flamingo-9B & Flamingo-80B\\
    \midrule
    冻结语言模型 & Chinchilla 1.4B；24层，宽2048，16头 & Chinchilla 7B（7.1B）；40层，宽4096，32头 & Chinchilla 70B；80层，宽8192，64头\\
    门控交叉注意力 & 每层插入；共24层，宽2048，16头；1.2B & 每4层插入；共10层，宽4096，32头；1.6B & 每7层插入；共12层，宽8192，64头；10B\\
    感知重采样器 & 6层，宽1536，16头，64个输出；0.194B & 同左 & 同左\\
    视觉编码器 & NFNet-F6，0.435B，冻结 & 同左 & 同左\\
    总参数 & 3.2B & 9.3B & 80B\\
    \bottomrule
  \end{tabularx}
  \caption{Flamingo三个规模的配置，取自原论文附录的Transformer超参数表和参数量表。只有门控交叉注意力层与感知重采样器参与训练。}
  \label{tab:vision-models-flamingo-config}
\end{table}

初始贡献为0不意味着门控学不到信号。若把后续误差对$\bm H'$的导数记为$\bm G$，则$\partial L/\partial\alpha=(1-\tanh^2\alpha)\langle\bm G,C(\bm H,\bm Z)\rangle$，其中内积对矩阵全部元素求和。$\alpha=0$时这个导数仍可非零；而该分支权重的直接梯度先乘以$\tanh\alpha$，在门控打开后才逐渐获得信号。教学上若原隐向量为$(1,0)^\top$、分支贡献为$(0,2)^\top$，门控从0变为0.2后，更新结果为$(1,0.3948)^\top$。

交错输入靠掩码控制可见范围。原Flamingo的每个文字位置直接交叉注意到它之前最近的一幅图像或一个视频，而不是同时直接读取全部历史图片；历史图像仍可通过此前文字隐状态和语言自注意力间接影响结果。示例中的红车和待答图中的红车虽然类别相同，也必须保持来源位置，不能把示例属性无条件移到目标上。

训练数据包含网页中的交错图文、图文对及视频文字对，目标仍是式\eqref{eq:vision-vlm-text-objective}的文字负对数概率。四个来源是约4300万网页的交错图文M3W、18亿对图像与替代文字ALIGN、3.12亿对长描述图文LTIP和2700万段短视频文字对VTP。对每个数据来源分别计算误差，再按来源权重（依次为1.0、0.2、0.2、0.03）累积梯度，以AdamW、学习率$10^{-4}$训练500K步，更新重采样器、新增门控层及相应新增词元；NFNet视觉编码器和Chinchilla语言模型的权重保持冻结。推理时把示例图$A$、示例答案“车窗关闭”、目标图$B$和问题依次放入上下文，目标答案位置直接读取$B$的64个视觉单元，再逐词解码；$A$的答案提供格式示范，却不能代替观察$B$的车窗。示例本身不引起参数更新，这就是少样本上下文学习；示例数、输入处理和上下文压缩仍限制能保存的细节。

\subsubsection{LLaVA}
\label{sec:vision-models-llava}

BLIP-2和Flamingo都设计了专门的压缩或读取模块，需要较多图文数据训练它们。LLaVA提出的问题是：如果语言模型本身足够强，能否只用最简单的连接，把训练重点放在“按指令回答图像问题”这一行为上？早期模型多在图像描述数据上训练，学会的是写一句描述，而不是按用户的多样请求组织回答。LLaVA的主要想法有两点：用一个线性投影把每个视觉网格位置直接映射到语言嵌入空间，不做压缩；再用视觉指令数据（对话、详细描述和复杂推理问答）训练语言模型按请求回答。本节依据原论文v2的线性投影设置，后来的多层感知机连接器不属于这里的模型\scite{vision-llava2023}。

视觉编码器是预训练的CLIP ViT-L/14，结构见表\ref{tab:vision-models-clip-config}。论文没有单列输入分辨率；按该编码器$224\times224$的训练分辨率，它输出$16\times16=256$个位置的网格特征$\bm F\in\mathbb R^{256\times1024}$，不含分类位置。论文比较了取最后一个Transformer层之前和之后的网格特征，在ScienceQA上前者略好（90.92\%对89.96\%），论文推测最后一层更偏向全局、抽象的图像性质，前一层保留更多局部性质。设$\bm W\in\mathbb R^{d_{\mathrm l}\times1024}$，投影得到$\bm P=\bm F\bm W^\top\in\mathbb R^{256\times d_{\mathrm l}}$，与问题词元的嵌入一起进入自回归语言模型Vicuna。视觉位置与文字位置之间的交互完全由语言模型自身的自注意力完成，和Flamingo另设交叉注意力分支不同。以两个二维教学网格$(1,0)$、$(0,1)$及$\bm W=\left[\begin{smallmatrix}2&0\\0&1\end{smallmatrix}\right]$为例，投影得到$(2,0)$、$(0,1)$。投影改变坐标宽度或尺度，但没有额外读取像素，也没有把每个图像块翻译成一个单词。语言模型主要是Vicuna-13B，论文也报告了7B版本；$d_{\mathrm l}$由所选Vicuna的嵌入宽度决定，论文未单列。整个连接器只有$\bm W$一个矩阵，参数量为$1024d_{\mathrm l}$，远少于BLIP-2的Q-Former或Flamingo的门控层。表\ref{tab:vision-models-llava-config}汇总各组件。

\begin{table}[htbp]
  \centering
  \small
  \begin{tabularx}{\linewidth}{@{}p{14mm}XXp{18mm}@{}}
    \toprule
    组件 & 输入与输出 & 结构与规模 & 训练状态\\
    \midrule
    视觉编码器 & 图像$\to256\times1024$网格特征 & CLIP ViT-L/14（24层，宽1024，16头） & 两阶段均冻结\\
    投影 & $256\times1024\to256\times d_{\mathrm l}$ & 一个线性矩阵$\bm W\in\mathbb R^{d_{\mathrm l}\times1024}$ & 两阶段均训练\\
    语言模型 & 视觉单元与文字嵌入$\to$回答词元 & Vicuna-13B（另有7B版本） & 第一阶段冻结，第二阶段训练\\
    \bottomrule
  \end{tabularx}
  \caption{LLaVA的组件与数据流，依据原论文v2。第一阶段用约59.5万对过滤后的CC3M图文训练1轮，学习率$2\times10^{-3}$，批次128；第二阶段用15.8万条指令样本训练3轮，学习率$2\times10^{-5}$，批次32；在8块A100上分别约需4小时和10小时。}
  \label{tab:vision-models-llava-config}
\end{table}

训练分为特征对齐和视觉指令适配两个阶段。对齐阶段用约60万个图文对构造“请简要描述图片”一类请求，冻结视觉编码器与语言模型，只更新投影层。冻结时$\bm F$可作为常量；若误差对$\bm P$的导数为$\bm G\in\mathbb R^{256\times d_{\mathrm l}}$，投影的梯度就是$\bm G^\top\bm F$，形状与$\bm W$相同，误差仍需经冻结的语言层传回投影输入。指令适配阶段继续冻结视觉编码器，同时更新投影与语言模型，数据为约15.8万条视觉指令样本，其中对话5.8万条、详细描述2.3万条、复杂推理7.7万条；对多轮数据，输入保留图片、用户请求和前轮回答，损失只监督助手应生成的回答词元。

指令数据的构造方式决定了模型学到什么。原论文用图片已有描述和对象框作为视觉内容的符号表示交给只读文字的教师模型，由它生成对话、详细描述和复杂问答；教师没有直接看原像素，但训练LLaVA时仍把真实图片交给视觉编码器。若请求是“比较两辆车的状态”，训练示例帮助模型学会“比较”这一请求形式，能否正确读取车的位置、朝向和遮挡仍由视觉证据及训练覆盖决定。若回答写了“右车正在行驶”，单张静态图片通常不足以确定实际速度；模型必须避免把语言中常见的情境补全当成观测。原模型可以生成自然语言，但没有因此得到坐标、实例掩码或图像生成能力。

\begin{figure*}[tbp]
  \centering
  % 每行仅保留用于区分接入位置的模块，不展开完整Transformer层。
\begin{tikzpicture}[x=1cm,y=1cm,font=\figurefont\small,
  block/.style={draw=FigureModel,fill=FigureModel!10,line width=0.8pt,
    minimum height=9mm,minimum width=21mm,align=center},
  input/.style={block,draw=FigureInput,fill=FigureInput!10},
  output/.style={block,draw=FigureOutput,fill=FigureOutput!10},
  arr/.style={-{Stealth[length=2mm]},BookInk,line width=0.8pt}]
  \node[anchor=west,font=\figurefont\small\bfseries] at (0,6.7) {BLIP-2：查询压缩后作为输入前缀};
  \node[input] (bf) at (1.1,5.7) {视觉特征};
  \node[block] (bq) at (4.1,5.7) {Q-Former};
  \node[block] (bp) at (7.2,5.7) {投影后的\\视觉前缀};
  \node[block] (bl) at (10.7,5.7) {语言模型};
  \node[output] (bo) at (14.1,5.7) {回答词元};
  \draw[arr] (bf) -- (bq);
  \draw[arr] (bq) -- (bp);
  \draw[arr] (bp) -- (bl);
  \draw[arr] (bl) -- (bo);
  \node[anchor=north] at (4.1,5.15) {可学习查询读取图像};
  \node (btext) at (10.7,4.7) {文字问题};
  \draw[arr] (btext) -- (bl);

  \node[anchor=west,font=\figurefont\small\bfseries] at (0,3.9) {Flamingo：在语言层之间读取视觉表示};
  \node[input] (ff) at (1.1,2.45) {视觉特征};
  \node[block] (fr) at (4.1,2.45) {Perceiver\\Resampler};
  \node[block] (fc) at (7.2,2.45) {门控交叉\\注意力};
  \node[block] (fl) at (10.7,2.45) {后续语言层};
  \node[output] (fo) at (14.1,2.45) {回答词元};
  \draw[arr] (ff) -- (fr);
  \draw[arr] (fr) -- node[above] {$K,V$} (fc);
  \draw[arr] (fc) -- (fl);
  \draw[arr] (fl) -- (fo);
  \node (fh) at (7.2,3.35) {语言隐状态$Q$};
  \draw[arr] (fh) -- (fc);
  \node[anchor=north] at (4.1,1.9) {每个图像／视频64个单元};

  \node[anchor=west,font=\figurefont\small\bfseries] at (0,0.7) {原始LLaVA：网格特征投影后进入共享序列};
  \node[input] (lf) at (1.1,-0.25) {视觉网格};
  \node[block] (lp) at (4.1,-0.25) {线性投影};
  \node[block] (ls) at (7.2,-0.25) {视觉单元\\与问题序列};
  \node[block] (ll) at (10.7,-0.25) {语言模型};
  \node[output] (lo) at (14.1,-0.25) {回答词元};
  \draw[arr] (lf) -- (lp);
  \draw[arr] (lp) -- (ls);
  \draw[arr] (ls) -- (ll);
  \draw[arr] (ll) -- (lo);
\end{tikzpicture}

  \caption{三种视觉接入方式的结构示意，按各原论文的预训练路径绘制。BLIP-2由少量查询形成视觉前缀，Flamingo在语言层中新增读取视觉表示的交叉注意力，原始LLaVA把视觉网格线性投影后交给语言主干。图中省略各模块内部层，框旁文字指出视觉信息进入的位置；冻结范围由正文所述阶段决定。}
  \label{fig:vision-language-connectors}
\end{figure*}
\FloatBarrier

\subsubsection{Qwen3-VL}
\label{sec:vision-models-qwen3vl}

LLaVA式的直接投影只在固定分辨率下读取视觉编码器最末层的特征，处理多页文档和长视频时出现两个不足。其一，最末层特征偏向高层语义，小字、细边界等低层细节在逐层抽象中可能减弱；其二，把视频帧当成一组没有时间含义的图片，模型无法知道两帧相隔多久。Qwen3-VL技术报告v1的主要想法是：视觉编码器接受动态分辨率输入，经空间合并后直接投影接入语言模型，同时把视觉编码器的多个中间层分别注入语言模型的前几层；视频则用位置编码与文字时间戳两种方式记录时间\scite{vision-qwen3vl2025}。

视觉编码器采用SigLIP-2结构，默认是SigLIP2-SO-400M，2B和4B两个小模型改用约3亿参数的SigLIP2-Large；它按原始宽高比接受动态分辨率输入，位置上使用二维旋转位置编码，并按输入尺寸插值绝对位置嵌入，输出长度随图像大小变化的视觉特征。语言主干有四个稠密规模（2B、4B、8B、32B）和两个混合专家规模（30B-A3B与235B-A22B，后者每个词元激活约220亿参数）；各规模的层数、宽度与头数，以及视觉编码器的块大小，技术报告未公开。表\ref{tab:vision-models-qwen3vl-config}按组件汇总。视觉语言合并模块使用两层\term{多层感知机}（multilayer perceptron, MLP），即在逐位置的线性映射之间加入非线性运算，把相邻$2\times2$个特征压缩为一个语言宽度的视觉单元，视觉单元数降为原来的四分之一。空间合并降低上下文长度，也把四个位置的信息汇集到一个表示中；实际细节保留需结合输入分辨率检查。

DeepStack多层接入利用视觉编码器三个不同深度的中间特征。每组特征经独立的合并模块映射后，以残差方式加到语言主干前三层相应视觉位置的隐状态上。这样，视觉信息从多个深度进入语言计算，而不是全部依赖初始输入；残差接入不增加序列位置，但增加了特征投影和存储。

位置表示有两项安排。交错多模态旋转位置编码（interleaved multimodal rotary position embedding, interleaved MRoPE）按时间、高度、宽度三个坐标旋转查询与键的部分分量，使注意力分数能够反映相对位置；它把三个坐标对应的旋转维度交错分散到不同频率范围，避免某个坐标只占据单一频段。视频时间戳则以文字形式放在每个时间块之前，例如“3.0秒”，提供明确的时间读数。在一个教学视频中，采样块对应1.0、3.0和8.0秒；保留时间戳，模型才有条件判断后两个观测相隔5秒。但如果敲门动作发生在4.2秒且没有进入任何采样块，时间戳再准确也无法补回动作。

训练从已有部件开始：视觉端由SigLIP-2检查点初始化，语言端由Qwen3初始化。预训练分四个阶段：S0只更新合并模块，冻结视觉编码器和语言主干，用约670亿词元、序列长8192让新接口先对齐；S1放开全部模块，用约1万亿词元做多模态预训练；S2把序列长度扩到32768，再训练约1万亿词元；S3以262144（即256K）的序列长度训练约1000亿词元，使模型适应长文档和长视频。

\begin{table}[htbp]
  \centering
  \small
  \begin{tabularx}{\linewidth}{@{}p{14mm}XXp{14mm}@{}}
    \toprule
    组件 & 输入与输出 & 结构与规模 & 训练状态\\
    \midrule
    视觉编码器 & 动态分辨率图像或视频帧$\to$变长视觉特征 & SigLIP2-SO-400M（2B、4B用SigLIP2-Large，约3亿参数）；二维旋转位置编码 & S0冻结，其后训练\\
    合并模块 & 相邻$2\times2$特征$\to$1个语言宽度单元 & 两层MLP & 各阶段训练\\
    DeepStack & 视觉编码器三个层级的特征$\to$语言前三层的残差 & 每个层级一个独立合并模块 & 同上\\
    语言主干 & 视觉单元与文字$\to$回答词元 & Qwen3：2B、4B、8B、32B稠密，30B-A3B与235B-A22B混合专家；层数与宽度报告未公开 & S0冻结，其后训练\\
    \bottomrule
  \end{tabularx}
  \caption{Qwen3-VL的组件，依据技术报告v1。DeepStack取用的具体视觉层号报告未给出。}
  \label{tab:vision-models-qwen3vl-config}
\end{table}各阶段样本都以图像或带时间戳的帧组为条件、以文字为目标，用式\eqref{eq:vision-vlm-text-objective}训练。后训练包含指令适配、蒸馏与强化学习，分别产生思考版和非思考版，二者响应行为不同，使用时应选定具体检查点。256K词元是容量约定，不是对任意长文档小字和短暂视频事件的保存保证。

\subsubsection{Qwen2.5-Omni}
\label{sec:vision-models-qwen-omni}

前面的模型只读取图像或视频，输出文字。视频中的敲门动作和敲击声可能互相补充，也可能错位或来自画外；用户还可能希望听到语音回答。这带来两个新问题：音频与视频怎样在时间上对应起来，以及怎样在生成文字的同时流式地生成语音，而不是等整段文字写完再交给另一个语音合成系统。Qwen2.5-Omni技术报告v1的主要想法是Thinker--Talker分工：Thinker是读取多模态输入、生成文字的语言模型，Talker读取Thinker的隐表示和已生成文字，生成语音编码，再恢复成声音\scite{vision-qwenomni2025}。

输入端，音频重采样到16kHz，按\ref{sec:vision-models-encoding}节的方法以25毫秒窗、10毫秒步长转为128频带的梅尔频谱，由Whisper-large-v3初始化的音频编码器形成时间特征，大约每40毫秒对应一个音频单元；流式场景下音频编码器的注意力按2秒一块计算，而不是覆盖整段录音。图像和视频由来自Qwen2.5-VL的视觉编码器处理，它是约6.75亿参数的ViT，块大小14，相邻$2\times2$个特征再由一个简单MLP合并为一个单元；单幅图像被当作两帧相同的视频，视频则按动态帧率采样，使帧时间与音频的时间刻度相配。时间对齐多模态旋转位置编码（time-aligned multimodal rotary position embedding, TMRoPE）给音视频提供共同时间尺度，一个时间位置单位对应40毫秒；音频单元沿时间排列，视频单元同时保留帧时间和空间位置。输入还按实际时间划分为2秒块，每块内先放视觉表示、后放音频表示，再按块交错。因此序列中相邻的视觉和声音单元由位置记录表达实际时间，而不是由先后次序表达。这个安排也不证明两段内容来自同一个声源：若门外另有人敲击金属，画面中的手部动作与声音可能在时间上接近，却缺乏确定来源。

输出端，Thinker按式\eqref{eq:vision-vlm-language}预测回答文字。Talker接收Thinker的高层隐表示以及已采样文字词元的嵌入，以双轨自回归方式预测离散语音编码；隐表示提供语义、语气等上下文，离散文字则帮助确定具体读音，避免语义接近但读音不同的内容仅由连续表示区分。语音编码还不是可以播放的波形：流式解码路径用流匹配模型把编码转为梅尔频谱，再由\term{声码器}（vocoder），即把频谱恢复为波形的网络，合成声音。流匹配的速度目标在\ref{sec:vision-models-flow-matching}节解释。解码器是用流匹配训练的DiT，注意力采用滑动窗口块掩码，每块的感受野共4块：回看2块、前看1块；声码器是修改过的BigVGAN，同样分块运行。因此流式输出分块产生，并非没有等待。

训练分工明确。Thinker的多模态预训练起步时固定语言模型，先训练连接和编码器，随后展开全参数与长上下文训练，再用多模态指令数据适配回答，目标都是文字负对数概率。Talker单独训练：以文字与对应朗读音频为记录，先用语音编码器把真实声音转为离散编码，再以此前编码、文字与Thinker隐表示预测下一个语音编码，目标编码的负对数概率提供监督；之后还有偏好优化和说话人适配。波形解码器以编码和真实频谱为配对学习频谱恢复，声码器学习从频谱到波形。理解监督不能替代语音发音监督，音画时间对齐也不能替代输出声学解码。表\ref{tab:vision-models-omni-config}汇总报告公开的组件信息。

\begin{table}[htbp]
  \centering
  \small
  \begin{tabularx}{\linewidth}{@{}p{14mm}XX@{}}
    \toprule
    组件 & 输入与输出 & 结构与规模\\
    \midrule
    音频编码器 & 16kHz音频、128频带梅尔频谱$\to$约每40毫秒一个单元 & Whisper-large-v3初始化；参数量报告未公开\\
    视觉编码器 & 图像或视频帧$\to$视觉单元 & 来自Qwen2.5-VL的ViT，约6.75亿参数，块14，$2\times2$合并\\
    Thinker & 交错的音视频与文字$\to$文字及隐表示 & Qwen2.5语言模型初始化，发布版本名为Qwen2.5-Omni-7B；层数与宽度报告未公开\\
    Talker & Thinker隐表示与已采样文字$\to$语音编码 & 双轨自回归Transformer解码器；规模报告未公开\\
    语音解码 & 语音编码$\to$梅尔频谱$\to$波形 & 流匹配DiT（滑动窗口块注意力，回看2块、前看1块）与修改版BigVGAN；规模报告未公开\\
    \bottomrule
  \end{tabularx}
  \caption{Qwen2.5-Omni的组件，依据技术报告v1。Thinker预训练第二阶段使用约8000亿图像与视频词元、3000亿音频词元和1000亿带声视频词元，前两阶段序列长8192，第三阶段扩到32768。}
  \label{tab:vision-models-omni-config}
\end{table}推理时，输入路径保留动作和声音，Thinker形成“画面中的人抬手敲门，并伴有敲击声”这样的回答，Talker把它组织为语音。该版本支持文字与语音输出，没有图像或视频生成模块。

\subsubsection{CounterCredit}
\label{sec:vision-models-observation}
\label{sec:vision-models-countercredit}

前面五个模型都只编码一次输入，然后连续生成回答。如果整图缩小后右车的车窗只剩几像素，语言模型继续生成更长的解释仍没有得到新的观测。另一种计算是暂停回答，指定原图中的车窗区域，由工具返回该区域的高分辨率像素，再继续语言预测。观察动作（visual observation action）指这次取回视觉内容的决定，不是答案中的一句“我再仔细看看”。它改变后续计算能够读取的信息，而仅延长原上下文中的文字生成不改变像素来源。

有了这个接口，还要学习什么时候值得调用。若只用最终是否答对作为强化学习回报，模型可能学会频繁调用工具，或者调用了却不使用返回内容；最终答对只能说明整条路径的结果，不能说明这次裁剪有必要，也不能说明答案依赖返回像素。2026年预印本CounterCredit的主要想法是，在已经发生的调用前状态上做两个对照：继续这条视觉分支与立即回答相比增加了多少正确答案得分，真实像素与同格式随机返回相比又增加了多少得分；只有两项都显著为正的调用才得到返还，其余调用付费。本节依据其v1方法、训练配置和返回格式，范围限于返回单幅裁剪图像的工具\scite{vision-countercredit-2026}。

模型架构仍是一个视觉语言模型，只是输出中可以出现调用动作。一次调用记录为$(h_i,c_i,o_i)$：$h_i$是调用前已有的完整状态，包括请求、先前图片与返回记录、已生成文字；$c_i$是此次动作及参数，即来源图片编号、坐标范围和输出尺寸；$o_i$是实际返回图像。工具取出像素后，模型重新编码$o_i$，把其表示放入相应图像轮次，随后才预测后续词元。若在先前裁剪中再次裁剪，应组合两次坐标变换才能回到原图。越界、空区域或调用格式无效时，没有可用的新观察，系统保留失败记录，而不是把失败当成已经看到了对象。

学习方法的关键是两个对照的得分。训练记录包含原图$\bm X$、问题$q$和参考答案$y$。策略$\pi_{\bm\theta}$从已有视觉语言检查点生成交错文字与裁剪的轨迹$\xi$，最后取得正确性回报$r_{\mathrm{out}}(\xi)\in\{0,1\}$。为了比较同一调用的不同返回，论文固定已实现的$h_i$：真实分支接入原动作$c_i$与原返回$o_i$；立即回答分支从$h_i$直接要求回答；随机分支固定$h_i$与$c_i$，只把返回像素替换成从原图均匀采样的同尺寸图块$\widetilde o_{i,k}$。随机图块不事先筛成“保证无关”，可能偶然含有答案证据，这是随机参照的限制。对于多选问题，设$O$为选项字母集合，$e$为固定的立即作答提示，参考答案得分为
\begin{equation}
  U_{\bm\theta}(p)
   =\frac{\pi_{\bm\theta}(y\mid p,e)}
          {\sum_{a\in O}\pi_{\bm\theta}(a\mid p,e)}.
  \label{eq:vision-countercredit-score}
\end{equation}
$p$是待探测的状态前缀，分母只在选项字母间归一化。得分不是最终抽样答对的频率，也不是外部判断器给整段推理打分。自由回答改用参考答案在该前缀下的每词元平均对数概率，单位为每词元奈特，与概率差的尺度不同。探测得分停止梯度，只用来计算回报。

令$U_i^{\mathrm{real}}=U_{\bm\theta}(h_i,c_i,o_i)$、$U_i^{\mathrm{now}}=U_{\bm\theta}(h_i)$，并用$K$个随机返回估计参照均值。两项差值为
\begin{equation}
  \begin{aligned}
    D_i&=U_i^{\mathrm{real}}-U_i^{\mathrm{now}},\\
    \overline U_i^{\mathrm{rand}}
      &=\frac1K\sum_{k=1}^{K}
          U_{\bm\theta}(h_i,c_i,\widetilde o_{i,k}),\\
    E_i&=U_i^{\mathrm{real}}-\overline U_i^{\mathrm{rand}}.
  \end{aligned}
  \label{eq:vision-countercredit-values}
\end{equation}
$D_i$称为决策价值（decision value），比较实际视觉分支与立即作答；$E_i$称为证据价值（evidence value），比较实际像素与同动作下的随机返回。前者还包含动作格式或延后回答的影响，后者尽量保持这些因素不变。若$D_i$大而$E_i$接近0，随机图块也带来同样收益，这次调用没有显示真实像素的额外作用；若$E_i$大而$D_i$接近0，像素相对于随机返回有信息，但立即作答已同样充分。因此只奖励其中一项会遗漏失败情形。论文取$a_i$为探测资格指示（有效单图返回、存在不同随机图块且得分有限），取$\varepsilon$为预先校准的误差死区，定义
\begin{equation}
  \begin{aligned}
    b_i&=\min\{E_i,D_i\}
        =U_i^{\mathrm{real}}
          -\max\{U_i^{\mathrm{now}},\overline U_i^{\mathrm{rand}}\},\\
    v_i&=r_{\mathrm{out}}(\xi)\,a_i\,\mathbb I[b_i>\varepsilon].
  \end{aligned}
  \label{eq:vision-countercredit-verification}
\end{equation}
$v_i=1$才表示此次调用获得双重核验。第二个等式来自同时减去两个参照时取较小差值，等价于比较较高的参照得分，所以任何一边没有超过死区都不能通过。$\varepsilon$用训练前的随机图块对照校准后固定，多选取0.05，自由回答取每词元0.44奈特，避免把有限$K$导致的小波动当成收益。

仍以车窗问题作教学计算。假设立即回答的正确选项得分为0.40，实际裁剪后为0.80，三个随机返回的得分为0.35、0.45、0.40，则$\overline U^{\mathrm{rand}}=0.40$，$D=E=b=0.40$，超过0.05，有效调用且最终答对时即可通过。若实际与随机返回都为0.80，则$D=0.40$而$E=0$，不能通过；若立即回答为0.82、实际为0.80、随机均值为0.40，则$E=0.40$而$D=-0.02$，也不能通过。三种情形展示了两个信号为何不能互相替代。

核验结果再转为回报。设$n(\xi)$为已执行调用数，$\mathcal V(\xi)=\{i:v_i=1\}$，核验通过的调用获得超过门槛部分的返还，未通过的调用付费：
\begin{equation}
  \begin{aligned}
    P(\xi)
      &=\min\left\{C,\,
         \gamma\sum_{i\in\mathcal V(\xi)}(b_i-\varepsilon)\right\}\\
      &\quad-\rho\bigl(n(\xi)-|\mathcal V(\xi)|\bigr),\\
    R(\xi)&=r_{\mathrm{out}}(\xi)+P(\xi).
  \end{aligned}
  \label{eq:vision-countercredit-price}
\end{equation}
$\gamma$把得分收益换成回报，$C$是每条轨迹返还上限，$\rho>0$是每个未核验调用的费用。若每条轨迹最多$n_{\max}$次调用且$\rho n_{\max}<1$，正确轨迹至少获得$1-\rho n_{\max}>0$，错误轨迹没有返还、总回报至多0，因此费用不改变正确与错误轨迹的排序。这只证明回报的界，不保证优化必然学会正确策略。用论文多选配置$\gamma=0.5$、$\rho=0.05$、$C=0.7$、$n_{\max}=6$复算，上述$b=0.40$的调用返还$0.5(0.40-0.05)=0.175$；若随后还有一个未通过调用，轨迹价格为0.125，答对时总回报为1.125。

策略更新使用\term{组相对策略优化}（group relative policy optimization, GRPO），即对同一个问题抽样一组轨迹，用组内相对回报决定哪些生成动作应提高概率；根据优势调整动作概率的原理可参考\BookChapterRef{chap:rl-learning}{卷七“学习”章}。CounterCredit把正确性和价格分开处理。若第$j$条轨迹的正确性为$O_j$、价格为$P_j$，组均值为$\overline O,\overline P$，$\sigma_{\min}=0.15$为标准差下限，其优势为
\begin{equation}
  \widehat A_j=
    \frac{O_j-\overline O}
         {\max\{\operatorname{std}(O),\sigma_{\min}\}}
       +(P_j-\overline P).
  \label{eq:vision-countercredit-advantage}
\end{equation}
正确性项除以组内标准差，价格项只减均值，所以多一次未核验调用造成的$\rho$差值保持原回报尺度。组内全错时，论文把价格优势设为0，避免只根据错误轨迹中的调用数学习。令$r_{j,k}$为轨迹$j$在生成位置$k$的新旧策略概率比，更新最大化$\min\{r_{j,k}\widehat A_j,\operatorname{clip}(r_{j,k},1-\delta,1+\delta)\widehat A_j\}$的位置平均，$\delta$限制概率比的有效更新幅度。参考答案、随机图块及探测得分不参与反向传播，梯度只来自已抽样动作的对数概率。

原实验从Qwen2.5-VL-7B-Instruct的裁剪轨迹监督微调开始，固定视觉编码器与投影器，随后在同一格式和调用预算下做三轮GRPO训练。推理时不提供参考答案、不运行对照探测，模型只根据当前状态选择立即回答或实际观察。训练中每个合格调用需要$K+2$次额外得分前向，这是训练成本。论文的信号不直接覆盖搜索、OCR文字返回或任意外部工具；随机参照偏差、策略得分误校准和分布变化也会限制核验的解释。

\begin{table*}[tbp]
  \centering
  \small
  \begin{tabularx}{\linewidth}{@{}p{22mm}p{30mm}p{30mm}XX@{}}
    \toprule
    模型 & 视觉进入语言计算的位置 & 每幅图的视觉单元 & 训练时更新的部分 & 新增的输入或输出\\
    \midrule
    BLIP-2 & 输入前缀 & 固定32个查询输出 & Q-Former与投影；视觉、语言冻结 & —\\
    Flamingo & 语言层间的门控交叉注意力 & 固定64个重采样输出 & 重采样器与新增层；视觉、语言冻结 & 图文交错上下文\\
    LLaVA & 输入序列中与文字并列 & 全部网格位置（256个） & 先投影，后投影与语言模型 & 视觉指令数据\\
    Qwen3-VL & 输入序列及前三层残差注入 & 随分辨率变化，$2\times2$合并 & 先合并模块，后全部模块 & 视频时间戳、长上下文\\
    Qwen2.5-Omni & 输入序列，音视频按时间交错 & 随分辨率和时长变化 & Thinker与Talker分别训练 & 音频输入、语音输出\\
    CounterCredit & 在已有模型上增加观察动作 & 按需追加裁剪图 & 策略整体，经监督微调与GRPO & 回答中途取回像素\\
    \bottomrule
  \end{tabularx}
  \caption{语言响应生成模型的比较。前三行均以单图问答为例；后三行在前者基础上扩展输入、输出或观察时机。}
  \label{tab:vision-models-language-compare}
\end{table*}

表\ref{tab:vision-models-language-compare}把六个模型放在同一单图问答任务下比较。BLIP-2与Flamingo都用固定数量的查询压缩视觉特征，换来的是语言计算量不随分辨率增长，代价是压缩时丢掉的细节无法再恢复；Flamingo把读取位置从输入前缀移到语言层之间，使多幅图与文字能按位置交错，代价是要在语言模型内部插入新层。LLaVA放弃压缩，保留全部网格位置，结构最简单，训练可以集中在指令数据上，代价是序列长度随视觉单元数增长。Qwen3-VL沿用LLaVA的直接投影，用$2\times2$合并控制长度，并以多层注入补回低层细节、以时间戳补回视频时间，代价是更长的上下文和更多的训练阶段。Qwen2.5-Omni把同一框架推广到音频输入和语音输出，新增的是时间对齐与独立的语音生成路径。CounterCredit不改连接结构，而是改变观察发生的时机：前五个模型的证据在一次编码时就已确定，CounterCredit允许回答过程中再取回像素，并用两个对照约束这种取回确有必要且被使用；代价是推理中的额外编码，以及训练中的额外探测前向。

\subsection{图像生成}
\label{sec:vision-models-media-generation}
\label{sec:vision-models-image-generation}

生成一幅雨天街景时，仅有“雨天”这个语义判断还不够，模型需要决定每处颜色、纹理与空间关系。输入是文字请求，有时还有参考图片；输出是一幅图像；难点有两个。其一是成本：$256\times256\times3$的图像已有约20万个数值，扩散模型要对这样的数组反复去噪数十到上百次。其二是控制：文字只描述了一部分内容，模型既要遵循文字，又要补出文字没有说明的细节。

扩散模型逐步加噪与去噪的原理已在卷六讲过：前向过程按噪声日程给数据逐步加入高斯噪声，网络学习预测所加噪声，生成时从纯噪声出发反复去噪，可参考\BookChapterRef{chap:pgm-generative}{卷六“概率生成模型”章}。图像生成主干（例如在潜空间中处理图像块序列的DiT）及图像压缩编码器（例如VQGAN）在\ref{sec:vision-ivm-generation}节讲过。本小节的三个模型都以扩散模型为基础，差别在于对上面两个难点的回答。LDM回答成本问题：先用自编码器把图像压缩为小得多的空间潜表示，再在潜空间中扩散，文字通过交叉注意力进入去噪网络。unCLIP回答“文字到像素一步跨得太大”的问题：先根据文字生成一个CLIP图像向量，再根据这个向量生成像素，把语义与细节分成两个阶段。GLIDE回答控制问题：在像素空间做文字条件扩散，并比较两种在采样时加强文字条件的引导方式。三者分别改变生成空间、中间表示和采样过程。

\subsubsection{LDM}
\label{sec:vision-models-ldm}

在像素空间训练扩散模型，大量计算花在人眼几乎察觉不到的高频细节上，训练和采样都要处理完整分辨率的数组。\term{潜空间扩散模型}（latent diffusion model, LDM）的主要想法是把“压缩”和“生成”分开：先训练一个自编码器，把图像压缩成保留可感知内容的低分辨率潜表示，并能高质量还原；再固定这个自编码器，只在潜空间中训练扩散模型。潜表示保留二维网格结构，因此去噪网络仍可使用卷积和空间注意力，这与CLIP把整图汇总成一个语义向量的空间不同\scite{pgm-generative-rombach2022}。

架构由三部分组成。图像自编码器的编码器$E$把$\bm X\in\mathbb R^{H\times W\times 3}$压缩为$\bm Z=E(\bm X)\in\mathbb R^{h'\times w'\times d_z}$，下采样因子$f=H/h'$，解码器$D$还原$D(\bm Z)\approx\bm X$；论文比较了多种$f$，以4到16之间较好地平衡了细节与成本。例如$f=8$、$d_z=4$时，$512\times512\times3$的图像变为$64\times64\times4$的潜网格，数值个数减少为约四十八分之一。去噪网络是带注意力层的U-Net，在潜网格上工作。条件编码器$E_{\bm\phi}$把文字$c$编码为词元序列$\bm C=E_{\bm\phi}(c)\in\mathbb R^{m\times d_c}$，并通过交叉注意力进入去噪网络：查询来自当前去噪网络的空间特征，键和值来自文字序列。若当前有$n_z$个空间单元、文字有$m$个词元，注意力权重是$n_z\times m$的矩阵，每个空间位置按当前状态读取不同文字信息；它没有直接为“红色”指定一个真实对象框，词语与空间的对应仍需通过数据和目标学习。

论文在LAION-400M上训练的文字到图像模型给出了一组具体配置，见表\ref{tab:vision-models-ldm-config}。它在$256\times256$图像上训练，使用$f=8$、4个潜通道的KL正则自编码器，潜网格为$32\times32\times4$。去噪U-Net的基础通道数为320，通道倍数为$(1,2,4,4)$，于是四个分辨率层级的潜网格边长依次为32、16、8、4，通道数依次为320、640、1280、1280；编码路径逐级下采样，解码路径对称地逐级上采样，并通过跳跃连接拼接同一层级的编码特征。表中每个层级的“深度”参数为2，论文未定义其精确含义。在潜网格边长为32、16、8的三个层级（不含最低的4）上，卷积残差块之后插入空间Transformer：先层归一化，用$1\times1$卷积改变通道，把$h'\times w'$个位置展平成序列，依次做自注意力、逐位置前馈和读取文字的交叉注意力，再用$1\times1$卷积变回特征图并加残差；这里Transformer的层数为1，注意力头数为8。文字编码器$E_{\bm\phi}$是一个从头训练的Transformer，用BERT分词器把提示切成至多77个词元，32层、宽1280，输出$\bm C\in\mathbb R^{77\times1280}$。整个模型约14.5亿参数。

\begin{table}[htbp]
  \centering
  \small
  \begin{tabularx}{\linewidth}{@{}p{14mm}XXp{14mm}@{}}
    \toprule
    组件 & 输入与输出 & 结构与规模 & 训练状态\\
    \midrule
    自编码器 & $256\times256\times3\leftrightarrow32\times32\times4$ & 下采样因子$f=8$，4个潜通道，KL正则 & 第一阶段训练，之后冻结\\
    去噪U-Net & $32\times32\times4$潜网格、噪声步、文字$\to$同形噪声预测 & 基础通道320，倍数$(1,2,4,4)$，层级边长32/16/8/4；在32、16、8层级插入1层空间Transformer，8头 & 第二阶段训练\\
    文字编码器 & 至多77个BERT词元$\to77\times1280$ & 32层Transformer，宽1280 & 与U-Net共同训练\\
    \bottomrule
  \end{tabularx}
  \caption{LDM文字到图像模型的配置，取自原论文附录的条件LDM超参数表与文字编码器表。扩散1000步，线性噪声日程，批次680，训练39万次迭代，学习率$10^{-4}$；整个模型约14.5亿参数。各层级的通道数由基础通道乘倍数算得。}
  \label{tab:vision-models-ldm-config}
\end{table}

学习分两个阶段。第一阶段训练自编码器，输入与目标都是同一张图：重建误差包括像素差异和已有感知网络特征的差异，局部判别器提高重建的真实程度，同时用很小权重的KL正则或向量量化约束潜空间，防止其尺度任意变化。这一阶段提高重建保真度，不训练“根据文字创造新场景”。第二阶段固定自编码器，在其潜空间训练扩散模型。把潜数组按固定次序展开为$\bm z$，取噪声步$k$，加入标准高斯噪声$\bm\epsilon$：
\[
  \bm z^{(k)}
  =\sqrt{\bar\alpha_k}\bm z+
    \sqrt{1-\bar\alpha_k}\bm\epsilon\text{，}
\]
其中$\bar\alpha_k=\prod_{j=1}^{k}\alpha_j$由噪声日程指定，$0<\alpha_j<1$，$(k)$表示加噪或去噪过程的步，不是视频帧。独立高斯加噪的合成仍为高斯，所以可以从干净潜表示一次构造任意$k$的带噪输入。去噪网络以当前潜表示、噪声步和文字为条件预测所加噪声：
\begin{equation}
  L_{\mathrm{LDM}}
  =\mathbb E_{\bm X,c,k,\bm\epsilon}
    \left[
      \left\lVert\bm\epsilon-
      \epsilon_{\bm\theta}
      \bigl(\bm z^{(k)},k,E_{\bm\phi}(c)\bigr)
      \right\rVert_2^2
    \right]\text{。}
  \label{eq:vision-ldm-objective}
\end{equation}
平方误差对输出的导数为$2(\widehat{\bm\epsilon}-\bm\epsilon)$，优化器把它传到去噪网络和文字编码器；原论文的文字条件设置共同优化$\bm\theta$与$\bm\phi$，后来使用冻结文字编码器的实现需另行说明。固定的自编码器提供目标坐标，不参与这次更新。

使用时，把带噪关系中的噪声替换成预测值，可以估计干净潜向量：
\begin{equation}
  \widehat{\bm z}
   =\frac{\bm z^{(k)}
       -\sqrt{1-\bar\alpha_k}\,
           \bm\epsilon_{\bm\theta}(\bm z^{(k)},k,\bm C)}
          {\sqrt{\bar\alpha_k}}.
  \label{eq:vision-ldm-clean-estimate}
\end{equation}
这来自对前向加噪等式移项，并非使用时已知真实$\bm z$。采用标准高斯反向采样时，记$\beta_k=1-\alpha_k$与$\widetilde\beta_k=\beta_k(1-\bar\alpha_{k-1})/(1-\bar\alpha_k)$，反向更新为
\begin{equation}
  \begin{aligned}
    \bm\mu_k
      &=\frac1{\sqrt{\alpha_k}}
        \left(\bm z^{(k)}
          -\frac{\beta_k}{\sqrt{1-\bar\alpha_k}}\,
             \bm\epsilon_{\bm\theta}(\bm z^{(k)},k,\bm C)\right),\\
    \bm z^{(k-1)}
      &=\bm\mu_k+\sqrt{\widetilde\beta_k}\,\bm\eta_k,
         \qquad \bm\eta_k\sim\mathcal N(\bm0,\bm I).
  \end{aligned}
  \label{eq:vision-ldm-reverse-interface}
\end{equation}
均值形式由前向高斯链的后验及噪声参数化得到，最后一步不再加随机噪声；不同采样器可以改用确定性或跳步更新。令$K$为噪声日程的末步，生成从近似标准高斯的$\bm z^{(K)}$开始，按$k=K,\ldots,1$反复预测，再把末态恢复为潜网格，用$D$解码成图片。

沿街景例子，只看潜网格中的一个两维单元：干净值为$(2,-1)^\top$，抽样噪声为$(0,1)^\top$，$\bar\alpha_k=0.64$，训练输入为$(1.6,-0.2)^\top$。若网络预测噪声$(0.1,0.8)^\top$，误差平方和为0.05，式\eqref{eq:vision-ldm-clean-estimate}恢复的干净估计为$(1.925,-0.85)^\top$，而不是原来的$(2,-1)^\top$。这条计算显示噪声误差怎样改变潜表示；它还不是完整街景像素，其他网格位置和解码器共同决定最终图像。压缩率每提高一倍，位置数按平方减少，但小字等细节也可能在压缩中丢失，生成后的事实与控制要求由\ref{chap:vision-image-generation}章的任务协议检查。

\subsubsection{unCLIP}
\label{sec:vision-models-unclip}

LDM让去噪网络直接从文字条件生成空间潜表示，一步跨越了“这段文字可能指哪些图像内容”和“这些内容具体长什么样”两层不确定性。另一方面，CLIP图像向量已经学到了与文字对齐的整图语义，并且对风格、构图等视觉属性也有编码。unCLIP的主要想法是在文字与像素之间插入CLIP图像向量这一中间表示：先验模型根据文字生成CLIP图像向量，解码器再根据该向量生成图片，从而把语义选择与像素细节分开处理。这是DALL-E 2报告中的路线，生成器并非CLIP本身\scite{vision-unclip2022}。

设文字为$c$，CLIP图像表示为$\bm z_{\mathrm i}$，图像为$\bm X$，模型写成
\begin{equation}
  p(\bm X\mid c)
  =\int p_{\mathrm{dec}}(\bm X\mid\bm z_{\mathrm i},c)\,
          p_{\mathrm{prior}}(\bm z_{\mathrm i}\mid c)\,
          \mathrm d\bm z_{\mathrm i}\text{。}
  \label{eq:vision-unclip-factorization}
\end{equation}
这里表达模型的分层生成分布，不宣称训练后已等于真实数据分布。先验处理“这段文字可能对应哪些图像语义表示”，解码器处理“这个语义表示可能对应哪些具体像素”；解码器还可以读取文字，不能把条件一律删成只有图像向量；训练时有10\%的样本把CLIP向量置零、50\%的样本丢弃文字，使解码器也能做无分类器引导。架构上，冻结的CLIP提供文字表示与图像表示；先验是一个Transformer；解码器是像素空间的扩散网络，结构改自GLIDE，图像表示被投影后加到噪声步嵌入中，并另投影为四个额外的上下文单元，与文字词元一起供注意力读取；之后两级扩散上采样模块把$64\times64$的图像依次放大到$256\times256$和$1024\times1024$，后两级以较低分辨率图像为条件。

论文附录给出了各部件的规模，见表\ref{tab:vision-models-unclip-config}。所用CLIP的图像编码器是在$256\times256$上训练的ViT-H/16，32层、宽1280；文字编码器24层、宽1024。解码器沿用GLIDE的35亿参数结构，在$64\times64$上生成，U-Net基础通道512，通道倍数$(1,2,3,4)$，每个注意力头64个通道，在32、16、8三个分辨率上做注意力。两个上采样器都不含注意力层，只用卷积残差块：$64\to256$的约7亿参数、基础通道320，$256\to1024$的约3亿参数、基础通道192。全卷积的上采样器只依赖局部邻域，避免了自注意力在高分辨率上随位置数平方增长的计算。两种先验都约10亿参数：扩散先验是宽2048、24个块的仅解码器Transformer，带因果掩码，输入序列依次为编码后的文字、CLIP文字向量、噪声步嵌入、加噪的CLIP图像向量，最后一个输出单元预测去噪后的图像向量；自回归先验的文字编码器宽2048、24块，解码器宽1664、24块、26个头，上下文长384。

\begin{table}[htbp]
  \centering
  \small
  \begin{tabularx}{\linewidth}{@{}p{16mm}XXp{12mm}@{}}
    \toprule
    组件 & 输入与输出 & 结构与规模 & 训练状态\\
    \midrule
    CLIP & 图像与文字$\to$表示向量 & 图像ViT-H/16（$256\times256$，32层，宽1280）；文字24层，宽1024 & 冻结\\
    扩散先验 & 文字、CLIP文字向量、加噪图像向量$\to$图像向量 & 仅解码器Transformer，24块，宽2048；约10亿参数；采样64步 & 训练\\
    解码器 & 图像向量与文字$\to64\times64$图像 & GLIDE式U-Net，基础通道512，倍数$(1,2,3,4)$，注意力于32/16/8；35亿参数；采样250步 & 训练\\
    上采样器一 & $64\to256$ & 无注意力的U-Net，基础通道320；约7亿参数；DDIM 27步 & 训练\\
    上采样器二 & $256\to1024$ & 无注意力的U-Net，基础通道192；约3亿参数；DDIM 15步 & 训练\\
    \bottomrule
  \end{tabularx}
  \caption{unCLIP各部件的配置，取自原论文附录的超参数表。CLIP在约6.5亿幅图像上训练，先验、解码器和上采样器只用约2.5亿幅图像的DALL-E数据集。}
  \label{tab:vision-models-unclip-config}
\end{table}

先验有两种训练路线。自回归先验先对训练图像表示作主成分分析，保留319个主成分，把每个分量的值分到1024个离散桶，再以文字及其CLIP表示为前缀逐分量预测桶编号，用负对数概率训练；使用时逐项抽桶、取对应数值、逆主成分变换，恢复候选图像表示。扩散先验保留连续向量：令$\bm z_{\mathrm i}^{(k)}$为按前述高斯规则加噪的图像表示，先验网络$f_{\bm\psi}$直接预测干净向量，
\begin{equation}
  L_{\mathrm{prior}}
   =\mathbb E_{\bm X,c,k,\bm\epsilon}
      \left[\left\lVert
         f_{\bm\psi}(\bm z_{\mathrm i}^{(k)},k,c)
           -\bm z_{\mathrm i}\right\rVert_2^2\right].
  \label{eq:vision-unclip-prior-objective}
\end{equation}
目标$\bm z_{\mathrm i}$由冻结CLIP计算，误差只更新$\bm\psi$。这里预测的是干净向量而非噪声，若采样器需要噪声形式，可用$\widehat{\bm\epsilon}=(\bm z_{\mathrm i}^{(k)}-\sqrt{\bar\alpha_k}\widehat{\bm z}_{\mathrm i})/\sqrt{1-\bar\alpha_k}$换算。解码器的训练以真实图像的CLIP表示和文字为条件，对带噪图像预测噪声，与式\eqref{eq:vision-ldm-objective}同类但作用在像素上。先验与解码器分别训练，CLIP始终冻结。

使用时，先验产生的$\widehat{\bm z}_{\mathrm i}$只表达一组整体条件。论文对每段文字采两个候选图像表示，按与CLIP文字表示的点积选其一；这是表示阶段的选择，尚未生成完整图片。固定这个表示、换不同图像噪声，解码器可以产生构图和纹理不同的多幅雨天街景，这正是分层的意义：相近的图像表示对应一族图片，表示匹配并不决定完整重建。参考图直接提供图像表示时可以跳过先验，得到保留主要语义的变体；但仅凭一个CLIP向量无法保留原图每处细节，要还原原图还需记录扩散末端状态等额外信息。

\subsubsection{GLIDE}
\label{sec:vision-models-glide}

条件扩散模型训练好以后，采样结果仍可能与文字只有模糊联系：网络学到的是$p(\bm X\mid c)$的近似，而许多训练描述并不精确，采样时条件的作用可能偏弱。在采样阶段加强条件，是独立于网络结构的另一个设计选择。GLIDE在像素空间训练文字条件扩散模型，并比较了两种引导方式：一种借助外部的图文匹配模型，用其梯度修正采样；另一种不用外部模型，用同一网络的条件与无条件预测之差外推。这为后续图像创作中的引导选择提供了具体入口\scite{vision-glide2022}。

架构上，基础模型约35亿参数，在$64\times64$分辨率上生成。其视觉部分是ADM结构的U-Net，即Dhariwal与Nichol为ImageNet $64\times64$设计的扩散网络，基础通道加宽到512，约23亿参数；文字部分是24个残差块、宽2048的Transformer，约12亿参数。文字Transformer把请求编码为$K$个词元，最后一个词元的表示替换ADM中原本的类别嵌入，与噪声步嵌入一起调制各残差块；整条$K$个词元的序列再投影到去噪网络每个注意力层的宽度，拼接到该层的注意力上下文中，作为额外的键和值，当前图像空间特征提供查询。这样文字既以整句摘要的形式影响全部卷积层，又以逐词元的形式被各空间位置按需读取。另一个约15亿参数的扩散上采样模型以文字和低分辨率图像为条件，把结果放大到$256\times256$；它的文字编码器宽1024，U-Net基础通道384。训练以所加噪声为目标，用平方误差训练文字编码器、去噪网络及连接权重，论文还学习反向过程的方差；基础模型以批次2048训练250万次迭代，上采样器以批次512训练160万次迭代。表\ref{tab:vision-models-glide-config}汇总这些配置。

\begin{table}[htbp]
  \centering
  \small
  \begin{tabularx}{\linewidth}{@{}p{16mm}XXp{12mm}@{}}
    \toprule
    组件 & 输入与输出 & 结构与规模 & 训练状态\\
    \midrule
    文字编码器 & 请求$\to K$个词元表示 & 24个残差块、宽2048的Transformer，约12亿参数 & 训练\\
    基础去噪网络 & $64\times64$带噪图像、噪声步、文字$\to$噪声与方差 & ADM式U-Net，基础通道512，约23亿参数（与文字编码器合计约35亿） & 训练\\
    上采样器 & $64\times64$图像与文字$\to256\times256$ & ADM上采样结构，基础通道384，文字编码器宽1024；约15亿参数 & 训练\\
    噪声适应CLIP & 带噪$64\times64$图像与文字$\to$匹配分数 & ViT-L CLIP，在加噪图像上重新训练 & 采样时冻结\\
    \bottomrule
  \end{tabularx}
  \caption{GLIDE的组件，取自原论文的训练细节。论文公开的过滤数据版本缩小为约3亿参数；U-Net的残差块数与注意力分辨率论文未给出。}
  \label{tab:vision-models-glide-config}
\end{table}

\term{无分类器引导}（classifier-free guidance, CFG）的想法是，让同一个网络既能做条件预测，又能做无条件预测，采样时沿两者之差加强条件。为此，训练时以一定概率（GLIDE取20\%）把文字替换为空序列，网络在这些样本上学到不读文字时的预测。令$\bm\epsilon_{\mathrm c}$与$\bm\epsilon_{\varnothing}$分别表示同一带噪图片、同一噪声步上的条件与空条件预测，采样时用
\begin{equation}
  \widehat{\bm\epsilon}
     =\bm\epsilon_{\varnothing}
        +g(\bm\epsilon_{\mathrm c}-\bm\epsilon_{\varnothing})
  \label{eq:vision-glide-cfg-interface}
\end{equation}
替代反向均值计算中的噪声预测，$g$是无量纲强度。$g=1$恢复条件预测，$g>1$沿条件与空条件预测的差向外推；这里的线性组合作用于噪声预测，不是把两幅完成图片混合。空条件预测是同一网络在真实训练过的空条件下运行，而不是事后删掉某个词向量。代价是每步需要两次网络前向。固定街景文字和一个二维教学噪声状态，假设$\bm\epsilon_{\varnothing}=(0.2,0.4)^\top$、$\bm\epsilon_{\mathrm c}=(0.1,0.8)^\top$，取$g=2$得到$(0,1.2)^\top$；第二维已经超出两种原预测的范围，说明引导是外推。增加强度可能提高条件遵循、降低样本多样性，也可能引入色彩饱和等异常，不能由组合形式推出越大越好。

CLIP引导需要另一套图文比较模型。它必须读取当前噪声水平的图像，所以GLIDE专门训练了适应带噪输入的CLIP，不能把只看清晰图片的原CLIP直接放在同一位置。设噪声适应模型给出匹配分数$s_{\bm\phi}(\bm x^{(k)},c)$，反向分布的原均值为$\bm\mu_k$、协方差为$\bm\Sigma_k$，均值修正为
\begin{equation}
  \widehat{\bm\mu}_k
   =\bm\mu_k+g_{\mathrm{clip}}\bm\Sigma_k
           \nabla_{\bm x^{(k)}}s_{\bm\phi}(\bm x^{(k)},c).
  \label{eq:vision-glide-clip-interface}
\end{equation}
梯度指出当前像素怎样局部变化会提高匹配得分，协方差把它转成与反向更新相容的均值修正。CLIP权重在采样中固定，求导对象是当前图像状态。教学上若$\bm\mu_k=(0.5,0.5)^\top$、$\bm\Sigma_k=0.1\bm I$、梯度为$(1,-2)^\top$、$g_{\mathrm{clip}}=2$，修正均值就是$(0.7,0.1)^\top$。这一修正发生在本步抽样之前，因而能改变未来轨迹；代价是每步额外的CLIP前向与反向。GLIDE的人工评价中，无分类器引导的样本在真实感和文字符合程度上更受偏好。

两种引导都与生成后重排不同。重排先完成$K$个候选，之后才计算匹配分数并选择；它不修改任何候选，也不能从全部错误的候选中创造正确答案。条件输入、采样引导和事后排序可以组合，但应把额外的网络调用、梯度和候选数量计入预算；图文相似度仍需与对象数量、属性归属和空间关系的独立核验配合。

\begin{table}[htbp]
  \centering
  \small
  \begin{tabularx}{\linewidth}{@{}p{18mm}XXX@{}}
    \toprule
    比较项 & LDM & unCLIP & GLIDE\\
    \midrule
    生成空间 & 自编码器的空间潜网格 & 先CLIP向量，后像素 & 像素\\
    文字条件进入 & 交叉注意力 & 先验读取文字；解码器读取图像向量与文字 & 噪声步嵌入与注意力上下文\\
    训练模块 & 自编码器先训后固定；去噪网络与文字编码器 & 先验、解码器、上采样器分别训练；CLIP冻结 & 文字编码器、去噪网络、上采样器\\
    采样时的控制 & 可配合引导 & 表示候选选择 & 无分类器引导或CLIP引导\\
    主要代价 & 压缩损失细节 & 两级生成，向量不保存细节 & 像素空间计算量大\\
    \bottomrule
  \end{tabularx}
  \caption{图像生成模型在文字到图像任务上的比较。}
  \label{tab:vision-models-image-compare}
\end{table}

表\ref{tab:vision-models-image-compare}固定文字到图像任务比较三者。LDM改变的是生成空间：把扩散移到压缩潜网格，换来训练和采样成本的大幅下降，付出的是自编码器重建损失，小字和细边界可能在压缩中丢失。unCLIP在文字与像素之间加入CLIP图像向量，换来同一语义下的多样变体和以图生图的能力，付出的是两级模型的训练成本，以及整图向量对数量、位置等细节的弱约束。GLIDE不改变生成空间，而在采样过程中加强条件；无分类器引导只需同一网络两次前向，CLIP引导需要额外训练噪声适应模型。三者可以组合：常见的文字到图像系统在LDM的潜空间中使用无分类器引导，unCLIP的解码器也使用了无分类器引导。

\subsection{视频与音频生成}
\label{sec:vision-models-flow-matching}
\label{sec:vision-models-video-audio-generation}

视频把图像生成的难点沿时间放大。输入是文字，有时还有首帧或已知片段；输出是数十到上百帧的画面，可能还有同步的声音。一段5秒、每秒16帧、$480\times832$的视频有约1.6亿个像素值，逐帧独立生成又无法保证同一物体在相邻帧中保持形状和位置。若还要配上声音，敲门的手接触门板的那一帧必须与敲击声对齐。因此需要回答三个问题：怎样把视频和音频压缩到可处理的规模，用什么目标训练生成主干，以及画面与声音怎样在生成中互相对应。

前两个问题有较一致的回答：用在时间和空间上同时压缩的变分自编码器得到潜表示，再用以图像块序列为输入的Transformer主干在潜空间生成，主干以流匹配目标训练。第三个问题区分了本小节的两个模型。Wan只生成视频，把文字、首帧和已知帧等条件统一接入一个视频主干；LTX-2让视频和音频各有一个分支，在生成的每一步通过交叉注意力交换信息，使声音与画面共同形成。两者都用流匹配训练，所以先说明这一目标。

媒体生成不一定把网络输出解释为“应去掉多少噪声”，也可以直接学习“当前表示应向哪个方向变化”。\term{流匹配}（flow matching）就是通过回归预先构造路径上的速度来训练连续生成模型。它要解决的问题是：用一条从噪声到数据的连续轨迹生成样本时，如何得到一个可以逐点监督的训练目标，而不必在训练中求解整条轨迹\scite{vision-flowmatching2023}。

用一个明确的教学构造说明。设数据潜向量为$\bm z_1\in\mathbb R^d$，噪声为$\bm z_0\sim\mathcal N(\bm0,\bm I)$，路径参数$s\in[0,1]$。在两者之间定义直线路径
\begin{equation}
  \bm z(s)=(1-s)\bm z_0+s\bm z_1,
  \qquad
  \bm u=\frac{\mathrm d\bm z(s)}{\mathrm ds}
       =\bm z_1-\bm z_0\text{。}
  \label{eq:vision-flow-path}
\end{equation}
$s=0$对应噪声，$s=1$对应数据，$s$是无量纲的生成进度，速度的单位是每单位生成进度的潜坐标变化，不是视频中的运动速度。这里采用无残余噪声的理想化端点以便手算；原流匹配论文也讨论终点保留微小方差等路径。

训练时取一对数据与条件$(\bm z_1,c)$，抽样噪声$\bm z_0$和时刻$s$，就能直接算出中间状态$\bm z(s)$与目标速度$\bm u$。速度网络$\bm v_{\bm\theta}$的平方误差目标为
\begin{equation}
  L_{\mathrm{FM}}
  =\mathbb E_{\bm z_1,c,\bm z_0,s}
     \left[
       \left\lVert
         \bm v_{\bm\theta}(\bm z(s),s,c)
           -(\bm z_1-\bm z_0)
       \right\rVert_2^2
     \right]\text{。}
  \label{eq:vision-flow-loss}
\end{equation}
计算一次训练目标不必先求解从0到$s$的微分方程，因为中间状态有显式构造。残差$\bm r=\bm v_{\bm\theta}-\bm u$的平方误差导数为$2\bm r$，再经速度网络反传到可训练权重；媒体编码器若固定，目标端点不因这次速度误差改变。流匹配提供目标构造，不给任意网络和数据配置一个自动收敛到真实分布的保证。

使用时却不知道目标$\bm z_1$，所以不能直接取$\bm z_1-\bm z_0$走向答案。模型从一个新噪声出发，根据当前状态预测速度，求解
\begin{equation}
  \frac{\mathrm d\widehat{\bm z}(s)}{\mathrm ds}
    =\bm v_{\bm\theta}(\widehat{\bm z}(s),s,c),
  \qquad
  \widehat{\bm z}(0)\sim\mathcal N(\bm0,\bm I)\text{，}
  \label{eq:vision-flow-ode}
\end{equation}
再把$\widehat{\bm z}(1)$送入媒体解码器。简单的欧拉法在$s_k$处用
$\widehat{\bm z}^{(k+1)}
=\widehat{\bm z}^{(k)}+(s_{k+1}-s_k)
\bm v_{\bm\theta}(\widehat{\bm z}^{(k)},s_k,c)$
更新。有限步长有数值误差，网络也有预测误差，两者不能混为同一项。

为什么从许多不同样本路径学习，仍能得到一个只依赖当前状态的速度？固定输入$(\bm z,s,c)$时，平方损失的最优预测是目标速度的条件均值。令$\bar{\bm u}=\mathbb E[\bm u\mid\bm z,s,c]$，则
\[
  \mathbb E[\lVert\bm v-\bm u\rVert_2^2\mid\bm z,s,c]
  =\lVert\bm v-\bar{\bm u}\rVert_2^2+
    \mathbb E[\lVert\bm u-\bar{\bm u}\rVert_2^2\mid\bm z,s,c]\text{。}
\]
展开平方后，交叉项因$\mathbb E[\bm u-\bar{\bm u}\mid\bm z,s,c]=\bm0$消失；右侧第二项不依赖$\bm v$，因此网络学习的是穿过当前状态的样本路径速度的条件平均。在速度场和密度足够规则时，这个平均速度场满足连续性方程$\partial_s p_s=-\nabla_{\bm z}\cdot(p_s\bar{\bm u})$，即局部流入流出的差决定该位置密度如何改变，从而推动整个分布沿路径从噪声变到数据。它不要求网络辨认训练时某条唯一的端点连线；有限数据和有限容量训练得到的$\bm v_{\bm\theta}$只是这一速度场的近似。

取二维教学样本$\bm z_0=(0,1)^\top$、$\bm z_1=(2,-1)^\top$，目标速度为$(2,-2)^\top$，$s=0.25$时状态为$(0.5,0.5)^\top$。假如使用时在该位置预测速度$(1.2,-2.8)^\top$，步长0.25给出下一状态$(0.8,-0.2)^\top$，而该训练直线在$s=0.5$处为$(1,0)^\top$。这个差异说明一次速度预测误差怎样传到位置，不能用单条训练直线宣称网络推理必然沿直线运行。图\ref{fig:vision-flow-step}给出几何对应。

\begin{figure}[htbp]
  \centering
  % 直线路径为教学构造；粗箭头使用特意偏离目标的速度预测。
\begin{tikzpicture}[x=3cm,y=1.65cm,font=\figurefont\small,
  arr/.style={-{Stealth[length=2mm]},BookInk,line width=0.8pt}]
  \draw[arr] (-0.15,-1.4) -- (2.55,-1.4) node[right] {$z_1$};
  \draw[arr] (-0.15,-1.4) -- (-0.15,1.5) node[above] {$z_2$};
  \foreach \x in {0,1,2}{
    \draw[BookMuted,line width=0.5pt] (\x,-1.4) -- (\x,-1.46);
    \node[below] at (\x,-1.46) {$\x$};
  }
  \foreach \y in {-1,0,1}{
    \draw[BookMuted,line width=0.5pt] (-0.15,\y) -- (-0.19,\y);
    \node[left] at (-0.19,\y) {$\y$};
  }
  \draw[BookBlue,line width=0.9pt] (0,1) -- (2,-1);
  \fill[BookBlue] (0,1) circle[radius=1.1mm];
  \fill[BookTeal] (2,-1) circle[radius=1.1mm];
  \node[anchor=south west] at (0,1.13) {$s=0,\ (0,1)$};
  \node[anchor=west] at (2.02,-0.96) {$s=1$};
  \fill[BookInk] (0.5,0.5) circle[radius=1mm];
  \node[anchor=south west] at (0.52,0.62) {$s=0.25,\ (0.5,0.5)$};
  \draw[BookTeal,line width=1pt,fill=white] (1,0) circle[radius=1.1mm];
  \node[anchor=west] (ideal) at (1.35,0.1) {路径上$(1,0)$};
  \draw[BookTeal,line width=0.6pt] (1.04,0) -- (1.3,0.1);
  \draw[-{Stealth[length=2.2mm]},BookAmber,line width=1.3pt]
    (0.5,0.5) -- (0.8,-0.2);
  \fill[BookAmber] (0.7825,-0.23) rectangle (0.8175,-0.17);
  \node[anchor=east] at (0.76,-0.35) {预测$(0.8,-0.2)$};
  \draw[BookMuted,dashed,line width=0.8pt] (0.82,-0.17) -- (0.97,-0.03);
  \node[anchor=west] at (0,-2.05) {粗箭头：$0.25(1.2,-2.8)^\top$};
\end{tikzpicture}

  \caption{流匹配教学样本的路径与一次欧拉更新。实线连接已知训练端点，粗箭头为预测速度乘以步长后的位移，虚线标示与训练路径上$s=0.5$位置的差别。横纵轴是两个潜变量坐标；路径参数$s$不是视频中的事件时间。}
  \label{fig:vision-flow-step}
\end{figure}

即便每条条件路径都是直线，整体数据分布对应的轨迹也未必是直线，更不能仅据此声称得到了全局最优传输。\term{规范化流}（normalizing flow）描述通过可逆变换建立分布的模型，流匹配描述训练速度场的一类目标，两者相关但不是同义词。若$\bm z_1=T(\bm z_0)$为可微可逆变换，密度须补上体积变化：
\begin{equation}
  \log p_1(\bm z_1)
    =\log p_0(\bm z_0)
       -\log\left|\det
          \frac{\partial T(\bm z_0)}{\partial\bm z_0}\right|.
  \label{eq:vision-flow-density-transform}
\end{equation}
例如$T(z)=2z$把同样概率质量铺在两倍长度上，所以密度减半。连续情形把许多微小变换相接；若速度场对空间可微且产生唯一可逆轨迹，沿轨迹的对数密度变化率为$-\nabla_{\bm z}\cdot\bm v_{\bm\theta}$，散度度量局部体积的扩张或收缩：
\begin{equation}
  \frac{\mathrm d}{\mathrm ds}
        \log p_s(\widehat{\bm z}(s))
    =-\nabla_{\bm z}\cdot
          \bm v_{\bm\theta}(\widehat{\bm z}(s),s,c).
  \label{eq:vision-flow-density-divergence}
\end{equation}
该式不是流匹配平方误差中的必选监督；只需生成样本的系统可以不计算似然，要求密度时才积分散度并计入求解代价。相关条件与变换原理见卷六“概率生成模型”章。

\subsubsection{Wan}
\label{sec:vision-models-wan}

在潜空间生成视频的直接做法，是把图像自编码器逐帧应用，再让生成主干处理逐帧潜表示的序列。这样时间维完全没有压缩，81帧就有81组潜网格，主干的序列长度和注意力成本都很高；而相邻帧高度相似，大量位置携带的是重复信息。另一方面，实际创作需要多种条件入口：只给文字、给定首帧、给定开头片段续写、给定首尾两帧补出中间。为每种入口单独训练一个模型代价很高。Wan技术报告v2的主要想法有两点：用因果三维变分自编码器在时间和空间上同时压缩视频，且第一帧单独编码，使图像可以当作单帧视频与视频共同训练；用一个以流匹配训练的Transformer主干在潜空间生成，文字经交叉注意力进入，首帧或已知帧则与噪声潜变量沿通道拼接，并用掩码说明哪些时间位置已知，从而用同一接口覆盖多种条件任务\scite{vision-wan2025}。

视频自编码器的核心是\term{因果卷积}（causal convolution），即时间方向上每个输出只依赖当前及此前帧的卷积，实现方式是只在时间轴的过去一侧填充。三维卷积把空间卷积沿时间扩展，因果约束使第一帧的编码不依赖后续帧，第一帧因此可以像一幅图像一样单独编码。设视频$\bm V\in\mathbb R^{(1+T)\times H\times W\times3}$，编码器把它压缩为潜表示$\bm z\in\mathbb R^{(1+T/4)\times(H/8)\times(W/8)\times16}$：第一帧只做空间压缩，其余帧每4帧压成一个潜时间位置，空间各压缩8倍，潜通道数为16。以$81$帧、$480\times832$的视频为例，潜表示为$21\times60\times104\times16$。这里的因果性属于自编码器，不表示生成主干只能从前向后逐帧预测。因果结构还使长视频可以分块编码：每块最多4帧，对应一个潜时间位置，块之间缓存前一块末尾的特征供下一块的卷积使用：时间核长为3的因果卷积缓存前一块的最后两帧特征，第一块则在前面补两帧零，内存因此不随视频总长增长。为保持时间因果，自编码器把会跨帧统计的组归一化换成逐位置的RMS归一化；空间上采样层的输入通道减半，推理显存约降低三分之一。整个自编码器约1.27亿参数。

生成主干沿用扩散Transformer的做法，其图像块化、噪声步调制等基础见\ref{sec:vision-ivm-dit}节。Wan用核为$(1,2,2)$的三维卷积把潜表示切成时空块，上例得到$21\times30\times52=32760$个视觉单元。文字由umT5编码为至多512个词元的嵌入序列，在每个Transformer块中通过交叉注意力被视觉单元读取。生成进度$s$先编码为向量，再由各块共享的一个小网络（线性层加SiLU激活）预测六个调制参数，用于缩放和平移各块中注意力和前馈层的输入输出，每块另学一组偏置。主干的输出与输入潜表示同形，是每个位置的预测速度。调制参数由所有块共享的网络统一预测，而不是每块各配一个，这显著减少参数；论文的消融把这一设置称为全共享，1.3B模型的30个注意力块都共享同一调制网络。

报告发布了1.3B和14B两个规模。报告正文给出的结构数字只有：1.3B模型有30个注意力块；两个规模的隐藏宽度、注意力头数、前馈宽度以及14B的层数，报告正文均未列出，这里不予补充。表\ref{tab:vision-models-wan-config}汇总报告可核实的组件信息。

\begin{table}[htbp]
  \centering
  \small
  \begin{tabularx}{\linewidth}{@{}p{14mm}XXp{12mm}@{}}
    \toprule
    组件 & 输入与输出 & 结构与规模 & 训练状态\\
    \midrule
    视频VAE & $(1+T)\times H\times W\times3\leftrightarrow(1+T/4)\times\frac H8\times\frac W8\times16$ & 三维因果卷积，RMS归一化，分块缓存；约1.27亿参数 & 先训练，之后冻结\\
    块化 & 潜表示$\to$时空单元序列 & 核$(1,2,2)$的三维卷积 & 训练\\
    文字编码器 & 提示$\to$至多512个词元嵌入 & umT5 & 冻结\\
    DiT主干 & 单元序列、生成进度、文字$\to$同形速度 & 1.3B：30个块，调制网络全共享；14B：层数、宽度与头数报告未列出 & 训练\\
    首帧分支 & 首帧$\to$全局上下文 & CLIP图像编码器加三层MLP，经独立交叉注意力接入 & 图生视频阶段训练\\
    \bottomrule
  \end{tabularx}
  \caption{Wan的组件，依据技术报告v2。图生视频模型的主干输入为$16+16+4=36$个通道。}
  \label{tab:vision-models-wan-config}
\end{table}

首帧和已知帧条件用拼接与掩码接入。给定首帧图像$\bm I$时，把它放在条件视频的第一个位置，其余帧填零，经同一个视频编码器得到条件潜表示$\bm z_{\mathrm c}$，形状与目标潜表示相同、16通道。零表示未提供内容，不是“未来画面应为黑色”的监督。另设掩码$\bm M$，约定1表示该帧已提供、0表示待生成，按原视频帧数逐帧记录，再按自编码器每4帧一组重排到通道维，得到4个通道，与潜时间位置对齐。含噪目标潜变量16通道、条件潜变量16通道和掩码4通道沿通道拼接，主干输入从16通道扩展为36通道；新增通道的输入投影零初始化，使模型起步时等同于原文字到视频模型。首帧还经CLIP图像编码器提取整体特征，通过三层MLP形成全局上下文，由与文字分开的另一组交叉注意力接入。条件潜表示保留空间位置，整体特征补充外观语义，两者不能互相替代。同一掩码接口只要改变哪些帧为1，就能表达续写（开头若干帧已知）、首尾帧转换（两端已知）和随机插帧（零散帧已知）。

学习分为自编码器与生成主干两部分。自编码器先作为二维图像自编码器训练，再把卷积扩展为三维因果卷积，在$128\times128$、5帧的短视频上训练，目标为L1重建误差、KL正则和感知误差的加权和，权重分别为3、$3\times10^{-6}$和3；最后在多种分辨率和帧数的高质量视频上微调，并加入三维判别器的对抗误差。主干按式\eqref{eq:vision-flow-loss}训练：对目标视频的干净潜表示$\bm z_1$和噪声$\bm z_0$构造$\bm z(s)=s\bm z_1+(1-s)\bm z_0$，回归速度$\bm z_1-\bm z_0$。生成进度$s$不按均匀分布抽取，而是取$s=\sigma(u)$，$u$服从正态分布，$\sigma$为logistic函数；这种logit正态抽样使中间进度被抽到的次数多于两端，训练更集中于从噪声到数据的中间阶段。训练先在$256\times256$图像上做文字到图像预训练，再进行图像与视频的联合训练，分辨率依次提高到192、480和720像素量级的视频，最后在480p和720p上后训练；优化器为AdamW，初始学习率$10^{-4}$。训练中只优化DiT主干，文字编码器和VAE编码器保持冻结，所以显存主要用于主干的激活与优化器状态。图像条件路线先在与文字到视频相同的数据上，用随机掩码混合训练各种条件任务，此时不接CLIP分支；随后对每种任务分别微调，并加入CLIP图像特征分支。

推理时，从与目标同形的噪声潜表示出发，在文字和条件下按式\eqref{eq:vision-flow-ode}求解，报告采用约50步并配合无分类器引导，最后由视频解码器分块恢复帧序列。固定同一首帧，分别要求玩具狗转头、向前走和镜头环绕，可以观察首帧条件留下了什么自由度：转头要补出原来看不到的侧面，行走要决定腿部运动和接触，镜头环绕还要补背景与背面。首帧只给了起点，这些未知内容不会因“图生视频”这个任务名称而变成已知。实际使用要对应具体检查点和条件入口，不能把报告中的多种任务并成每个权重都支持的功能。

\subsubsection{LTX-2}
\label{sec:vision-models-ltx2}

Wan只生成画面。若先生成敲门画面，再由另一个模型补声音，后一步只能适应已固定的动作，画面也无法因声音节奏而调整；反过来先生成声音再配画面也有同样问题。若把音频和视频压进同一个潜空间、用一个主干处理，二者的时间分辨率、通道含义和信息量差别很大，会互相挤占容量。LTX-2技术报告v1的主要想法是让音频和视频各有一个生成分支，分别使用适合自己的潜空间，同时在每一层通过双向交叉注意力交换信息，使两种媒体在同一个生成过程中共同形成\scite{vision-ltx22026}。

视频经因果\term{变分自编码器}（variational autoencoder, VAE）编码为空间与时间潜网格，音频经梅尔频谱转换和独立的音频VAE编码为时间潜序列。VAE把媒体压缩到受分布约束的连续表示，再由对应解码器还原，其原理见卷六“概率生成模型”章。两个分支的长度、通道宽度和压缩方式可以不同，因此不需要把音频潜向量解释成视频像素。每个分支有自己的自注意力和前馈层；视频分支内部使用时空位置，跨分支连接则只使用时间位置，因为音频单元没有图像平面上的横纵坐标。视频查询读取音频的键和值，音频查询反过来读取视频的键和值。文字条件来自冻结的Gemma3语言模型的多层隐表示，经特征提取投影及每个模态各自的文字连接器处理，再分别进入两个分支。时间编码使对应关系可被表达，但正确的敲击同步仍要从配对数据中学习、从生成结果中检验。

两个分支的容量是不对称的。报告摘要给出视频分支约140亿参数、音频分支约50亿参数，合计约190亿；结论部分却写成在130亿参数的视频扩散Transformer上扩展一个30亿参数的音频分支，两处口径不一致，报告未作解释，这里以摘要和推理性能一节的“19B”为准。不对称的理由是画面合成比声音合成需要更多容量。每个双分支块包含本分支的自注意力、读取文字的交叉注意力、音视频双向交叉注意力和前馈层，各操作之间穿插RMS归一化；生成进度通过跨模态共享的自适应层归一化同时调制两个分支。位置编码上，视频用覆盖横、纵、时间三个坐标的旋转位置编码，音频只用时间一维的旋转位置编码，跨分支注意力只取时间分量。

音频路径的数字较为具体：输入转为16kHz立体声，两个声道各算一张梅尔频谱，沿通道拼接后送入因果音频VAE；输出的每个潜单元约对应$1/25$秒，是一个128维向量，即每秒约25个音频单元。声码器基于HiFi-GAN，生成器通道数比HiFi-GAN V1加倍，以24kHz输出双声道波形。文字端以Gemma3-12B为主干，特征提取器取全部解码层的输出，经均值中心化缩放后展平，再用可学习矩阵投影；每个分支各有一个文字连接器，由若干双向注意力的Transformer块构成，并在序列中加入可学习的“思考”词元。视频VAE的压缩倍数、两个分支的层数、宽度和头数，以及连接器的块数，报告均未公开。表\ref{tab:vision-models-ltx2-config}按组件汇总。

\begin{table}[htbp]
  \centering
  \small
  \begin{tabularx}{\linewidth}{@{}p{14mm}XXp{14mm}@{}}
    \toprule
    组件 & 输入与输出 & 结构与规模 & 训练状态\\
    \midrule
    视频VAE & 视频$\leftrightarrow$时空潜网格 & 时空因果VAE；压缩倍数报告未公开 & 冻结\\
    音频VAE & 16kHz立体声梅尔频谱$\leftrightarrow$每秒约25个128维单元 & 因果VAE & 冻结\\
    文字编码 & 提示$\to$两个分支各自的条件序列 & Gemma3-12B（冻结）、多层特征提取投影、每分支一个带思考词元的连接器 & 投影先训后固定；连接器训练\\
    视频分支 & 视频潜单元$\to$速度 & 约140亿参数；层数与宽度报告未公开 & 训练\\
    音频分支 & 音频潜单元$\to$速度 & 约50亿参数；层数与宽度报告未公开 & 训练\\
    声码器 & 梅尔频谱$\to$24kHz双声道波形 & HiFi-GAN式，生成器通道加倍 & 冻结\\
    \bottomrule
  \end{tabularx}
  \caption{LTX-2的组件，依据技术报告v1。分支参数取摘要口径；报告称可连续生成至多20秒的音视频。}
  \label{tab:vision-models-ltx2-config}
\end{table}

训练把配对音画编码为各自的目标潜表示，构造噪声状态，用流匹配分别监督两个分支的速度。可以把联合状态写成$(\bm z^{\mathrm v}(s),\bm z^{\mathrm a}(s))$，两边的速度预测都依赖当前音画状态及文字条件。公开训练源码给出了符号方向：它以$\sigma=1$表示噪声、$\sigma=0$表示干净数据，构造$(1-\sigma)\bm z_{\mathrm{clean}}+\sigma\bm\epsilon$，目标为$\bm\epsilon-\bm z_{\mathrm{clean}}$。令$s=1-\sigma$，便恢复式\eqref{eq:vision-flow-path}的噪声到数据方向，同时速度需变号；不能在减小$\sigma$的采样器中误用本章的正向速度。源码中的音视频批次共用$\sigma$，各自抽同形噪声，各自按有效元素平均平方误差，最后把两项相加；作为条件的已知视频位置可由掩码排除监督。报告的初期阶段共同训练特征投影与生成模型，之后固定该投影，文字连接器仍与生成模块共同训练；两个媒体自编码器和声码器提供固定的解码空间。这个可核查实现不代表已公开全部基础预训练日程。

若把一个敲门教学片段的状态简化为视频和音频各一个标量，噪声端为$(0,1)$、数据端为$(2,-1)$，目标正向速度为$(2,-2)$。$s=0.25$时联合状态为$(0.5,0.5)$；若两分支预测为$(1.5,-1.5)$，总平方残差为0.5，步长0.25更新到$(0.875,0.125)$。真实系统有长时间潜序列与双向读取，例子只显示两种媒体在同一次生成步中各有目标、误差和更新。使用时，从两个潜空间的噪声开始共同更新；视频潜表示由视频VAE解码为帧，音频潜表示先恢复梅尔频谱，再由声码器恢复波形。若文字要求“人物敲门两次，随后停下”，验收至少要对应两次接触、两次声音以及停顿。LTX-2中的文字编码器服务生成条件，系统没有自由问答输出头。

\begin{table}[htbp]
  \centering
  \small
  \begin{tabularx}{\linewidth}{@{}p{18mm}XX@{}}
    \toprule
    比较项 & Wan & LTX-2\\
    \midrule
    输出 & 视频 & 视频与同步音频\\
    潜空间 & 一个视频因果VAE & 视频因果VAE与音频VAE各一\\
    主干 & 单一Transformer & 视频、音频双分支，层间双向交叉注意力\\
    条件入口 & 文字交叉注意力；已知帧经通道拼接与掩码；首帧CLIP特征 & 文字经各模态连接器进入两个分支\\
    训练目标 & 视频潜表示的流匹配速度 & 两分支速度误差之和\\
    \bottomrule
  \end{tabularx}
  \caption{视频与音频生成模型在文字到媒体任务上的比较。}
  \label{tab:vision-models-video-audio-compare}
\end{table}

表\ref{tab:vision-models-video-audio-compare}把两个模型的比较固定在文字到媒体的生成任务上。Wan与LTX-2都使用因果VAE压缩视频、在潜空间用Transformer主干和流匹配生成，差别在于LTX-2增加了一个音频分支和两分支之间的双向交叉注意力。这一改变换来同一生成过程中的音画对应，画面与声音能互相调整；代价是两套潜空间与解码器、跨分支注意力的额外计算，以及对配对音画训练数据的依赖。Wan则把设计重点放在单一视频主干的多种条件入口上，用通道拼接和掩码统一了图生视频、续写和首尾转换，但不输出声音，配声需要另接模型，\ref{chap:vision-multimodal-generation}章再讨论这两种组织的取舍。

\section{理解与生成的结合}
\label{sec:vision-models-unified}

前面两节的理解模型与生成模型各自成立。若用户希望“解释这张图，再按要求画一个变体”，或者希望生成的图片在交付前被检查并改正，就需要让理解与生成配合。问题在于两者之间传递什么、由谁训练：理解模型读到的证据怎样交给生成器，生成的结果又能否回到理解模型被检查。判断“结合”的程度，需要看哪些参数、表示和训练目标实际共用，不能只看界面是否只有一个输入框。

本节按结合的紧密程度排列三种做法，因为紧密程度决定了信息能否在两类能力之间无损传递，也决定了训练与替换模块的难度。最松的是基于接口调用的结合：理解模型和生成器各自独立训练，理解模型把内容说明写成文字或区域交给生成器，信息只能经这一中间格式传递，但两边都可以单独替换。更紧的是基于共享表示的结合：同一主干读取共同的视觉表示，用不同输出头分别产生文字和图像，两种能力在训练中共享参数。最紧的是基于反思反馈的结合：在共享主干之上，生成的图片再被同一模型读取和评判，评判结果又作为编辑条件生成修订图，理解与生成形成闭环。

\subsection{基于接口调用的结合}
\label{sec:vision-models-composition}

最直接的做法是不训练任何新模型，由视觉语言模型读取图片和请求，形成明确的内容说明，再调用专用生成器。它要解决的问题是：已有的理解模型和生成器各自很强，能否不经联合训练就完成需要两种能力的请求。以“把玩具狗旁的花瓶改成蓝色，其余尽量保持”为教学请求，理解模型（例如Qwen3-VL）可以输出目标短语“花瓶”、保留对象“玩具狗”和修改属性“蓝色”；若生成器支持局部编辑，还需要用定位模型把目标短语转换成原图坐标中的掩码，再把原图、掩码和改写后的文字一起交给编辑模型。

这条流水线的架构就是接口格式。传什么决定了信息损失：只传一句“蓝色花瓶旁有玩具狗”时，生成器没有收到原图的纹理和位置，只能依据语义重建场景；传原图与掩码可以保留更多空间条件，但若理解模型把目标区域选成玩具狗，生成器仍可能准确执行了错误区域上的请求。连续特征也可以通过适配模块传给生成器，减少必须把全部细节说成文字的限制；然而不同编码器的特征坐标和尺度未必兼容，适配模块需要自己的训练数据。因此，这种结合有三种训练安排：分别训练后直接连接、固定两端只学习适配器、共同更新多个模块，不能仅因推理时连续调用就称为共同学习。组合系统应保存中间说明、区域来源、输入图片和最终结果，使错误能定位到读取、传递还是生成环节。

\subsection{基于共享表示的结合}
\label{sec:vision-models-showo2}

接口调用的不足在于两边的表示互不相通：理解模型看到的细节只能通过文字或掩码转述给生成器，生成器也无法利用理解模型学到的语义。若让同一个主干同时负责理解和生成，就必须解决一个矛盾：理解需要高层语义特征，生成需要纹理和局部结构，而以往的统一模型往往只用其中一种视觉表示。Show-o2的主要想法是在同一视觉潜空间中提取这两类信息并融合为共享视觉表示，送入共享语言主干，再用两个输出头分别处理文字和视觉：语言头预测下一词元，流头用流匹配预测视觉潜表示的速度。本节依据论文v3\scite{vision-showo22025}。

图像和视频先由Wan2.1的三维因果VAE编码为潜网格，图像作为单帧情形处理。潜表示经过两条路径：语义层从潜网格中提取高层信息，它由预训练SigLIP的视觉层加新的潜块嵌入组成，并预先用SigLIP在干净图片上的特征蒸馏，使其能从干净和带噪潜表示中提取语义；投影层保留低层特征。设语义路径输出为$\bm S\in\mathbb R^{n_z\times d_{\mathrm s}}$、低层路径输出为$\bm P\in\mathbb R^{n_z\times d_{\mathrm p}}$，同一行须来自相同潜网格范围，沿特征维拼接为$[\bm S,\bm P]\in\mathbb R^{n_z\times(d_{\mathrm s}+d_{\mathrm p})}$，再经均方根归一化和两层映射得到语言宽度的共享视觉表示$\bm U$。这里的共享表示保留多个位置，并非把整幅图压成一个CLIP向量。干净参考图以路径终点$s=1$的状态进入，生成中的视觉块则带有当前路径时刻，所以同一视觉接入模块需要处理不同噪声水平。

用论文第一阶段的输入尺寸走一遍形状。$432\times432$的图像经Wan2.1 VAE在空间上压缩8倍，得到$54\times54$的潜网格；语义层和投影层都用$2\times2$的块嵌入，各得$27\times27=729$个位置，恰与SigLIP在$384\times384$输入、块大小14时的输出位置数相同，所以语义层可以直接以SigLIP的特征为蒸馏目标。投影层只是一个二维块嵌入层；语义层复用SigLIP-so400m的视觉Transformer块，换上新的$2\times2$块嵌入，预先以批次512蒸馏20万次迭代，与SigLIP特征的平均余弦相似度约0.9。两路特征拼接后经RMS归一化和两层MLP映射到语言宽度，并在序列前放一个生成进度的嵌入，干净图像取1.0。其他分辨率通过双三次插值位置嵌入适配。语义层的层数、流头的层数，以及模型的总参数量，论文未给出；表\ref{tab:vision-models-showo2-config}汇总可核实的部分。

\begin{table}[htbp]
  \centering
  \small
  \begin{tabularx}{\linewidth}{@{}p{14mm}XXp{14mm}@{}}
    \toprule
    组件 & 输入与输出 & 结构与规模 & 训练状态\\
    \midrule
    VAE & $432\times432$图像$\leftrightarrow54\times54$潜网格；视频时间压缩4倍 & Wan2.1三维因果VAE，空间8倍 & 始终冻结\\
    语义层 & 潜网格$\to729\times d_{\mathrm s}$ & SigLIP-so400m视觉块加新的$2\times2$块嵌入，经蒸馏；层数论文未给出 & 第二阶段训练\\
    投影层 & 潜网格$\to729\times d_{\mathrm p}$ & 一个$2\times2$块嵌入层 & 两阶段训练\\
    融合 & $[\bm S,\bm P]\to729$个语言宽度单元 & RMS归一化加两层MLP & 两阶段训练\\
    语言主干 & 交错的文字与视觉单元$\to$隐状态 & Qwen2.5-1.5B-Instruct或7B-Instruct & 第二阶段训练\\
    流头 & 隐状态$\to$潜网格同形速度 & 若干带adaLN-Zero调制的Transformer层；层数论文未给出 & 两阶段训练\\
    \bottomrule
  \end{tabularx}
  \caption{Show-o2的组件，依据论文v3的1.5B设置。第一阶段只训练投影层、融合模块与流头，第二阶段训练除VAE外的全部模块。}
  \label{tab:vision-models-showo2-config}
\end{table}

共享主干来自Qwen2.5-Instruct。文字嵌入和视觉块按交错顺序进入主干；序列块之间遵守因果顺序，使后续输出可以依赖先前条件；同一视觉块内部允许完整注意力，使待生成块中的各位置共同参与当前速度预测。语言头把隐状态转换为下一词元分布；流头包含带生成时刻调制的Transformer层，输出与视觉潜网格同形的速度，经数值更新和VAE解码才成为媒体。比如一个$27\times27$的视觉网格有729个空间单元，流头对各潜块预测若干连续分量，语言头则在一个回答位置对每个词元编号计分；这一差别解释了“统一”为什么不要求相同的输出头。图\ref{fig:vision-showo-heads}把共享范围与两个输出头分开表示。

\begin{figure*}[tbp]
  \centering
  \begin{tikzpicture}[x=1cm,y=1cm,font=\figurefont\small,
  block/.style={draw=FigureModel,fill=FigureModel!10,line width=0.8pt,
    minimum height=9mm,minimum width=20mm,align=center},
  input/.style={block,draw=FigureInput,fill=FigureInput!10},
  output/.style={block,draw=FigureOutput,fill=FigureOutput!10},
  arr/.style={-{Stealth[length=2mm]},BookInk,line width=0.8pt,
    rounded corners=1.2mm}]
  \node[input] (z) at (1.1,1.8) {干净／带噪\\视觉潜表示};
  \node[block] (sem) at (4.0,2.7) {语义层};
  \node[block] (low) at (4.0,0.9) {低层投影};
  \node[block] (fuse) at (6.9,1.8) {按位置拼接\\与融合};
  \node[block] (trunk) at (9.8,1.8) {共享主干};
  \node[block] (text) at (12.7,2.7) {语言头};
  \node[block] (flow) at (12.7,0.9) {流头};
  \node[output] (words) at (15.4,2.7) {回答词元};
  \node[output] (media) at (15.4,0.9) {更新潜状态\\VAE解码};
  \draw[BookInk,line width=0.8pt] (z.east) -- (2.55,1.8);
  \draw[arr] (2.55,1.8) |- (sem.west);
  \draw[arr] (2.55,1.8) |- (low.west);
  \draw[arr] (sem.east) -- (5.45,2.7) |- (fuse.west);
  \draw[arr] (low.east) -- (5.45,0.9) |- (fuse.west);
  \draw[arr] (fuse) -- (trunk);
  \draw[BookInk,line width=0.8pt] (trunk.east) -- (11.25,1.8);
  \draw[arr] (11.25,1.8) |- (text.west);
  \draw[arr] (11.25,1.8) |- (flow.west);
  \draw[arr] (text) -- (words);
  \draw[arr] (flow) -- (media);
  \node (prompt) at (9.8,0.05) {文字上下文};
  \draw[arr] (prompt) -- (trunk);
  \node[anchor=west] at (0,-0.75) {视觉块之间保持因果顺序，块内位置允许完整注意力。};
\end{tikzpicture}

  \caption{Show-o2的共享视觉接入与双输出头示意，依据论文v3重绘。视觉潜表示的语义路径和低层路径按位置融合，再与文字一起进入共享主干。语言头给出词元，流头给出视觉潜变量的速度；生成分支需迭代更新潜状态并重新经过模型，图中省略迭代回边以突出共享和专用模块的边界。}
  \label{fig:vision-showo-heads}
\end{figure*}
\FloatBarrier

学习目标把两种误差相加：
\begin{equation}
  L=\lambda_{\mathrm{text}}L_{\mathrm{NTP}}+L_{\mathrm{FM}}\text{，}
  \label{eq:vision-showo-loss}
\end{equation}
其中$L_{\mathrm{NTP}}$是文字的下一词元预测损失，即式\eqref{eq:vision-vlm-text-objective}；$L_{\mathrm{FM}}$是式\eqref{eq:vision-flow-loss}的视觉速度误差；$\lambda_{\mathrm{text}}$控制文字项权重。理解样本把干净图像块和问题作为条件，监督答案文字；生成样本把文字描述作为条件，对目标图像潜表示抽噪声和路径时刻，监督流头速度；交错样本让先前图像与文字成为后来图像的条件，不对作为条件的图片施加文字目标。对共享权重，一次联合更新的梯度是$\lambda_{\mathrm{text}}\nabla L_{\mathrm{NTP}}+\nabla L_{\mathrm{FM}}$；两个梯度可能指向不同方向，共享可以传递信息，也可能产生任务冲突，不能仅据共享结构保证两项能力同时改善。训练分两阶段：第一阶段只训练投影、时空融合和流头，保留已有语言知识，学习视觉生成接口；第二阶段用高质量理解与生成数据适配除VAE外的整体模型。VAE始终不更新。原论文1.5B设置以AdamW训练：第一阶段在约6600万图文对上以常数学习率$10^{-4}$训练15万次迭代，$\lambda_{\mathrm{text}}=0.2$，再在1600万条筛选后的生成数据上续训4万次迭代；第二阶段在约900万条多模态指令与1600万条生成数据上训练约3.5万次迭代，$\lambda_{\mathrm{text}}=1.0$。训练时以10\%的概率丢弃描述，使流头能做无分类器引导；视频与交错版本另外加入相应样本，较大的7B变体没有沿用同样的视频生成训练数据。

对于“解释参考图，再生成雨天变体”，参考图先按$s=1$送入，语言头逐词给出“右侧是红色汽车”；生成开始标记出现后，系统分配新的视觉噪声块，保留此前图片和文字作为条件。假设新块的一个二维单元在$s=0.25$为$(0.5,0.5)^\top$，流头预测$(1.2,-2.8)^\top$，步长0.25更新为$(0.8,-0.2)^\top$；更新后的块必须重新经过双路径融合、共享主干与流头，才得到下一时刻速度。完整视觉块到$s=1$后经VAE解码成图片。解释正确并不自动保证生成正确，生成中仍可能丢失刚才已识别的对象关系；该模型也没有音频生成能力。

\subsection{基于反思反馈的结合}
\label{sec:vision-models-unireason}

共享表示让理解和生成共用参数，但两者在一次请求中仍是单向的：先理解，后生成，生成结果不再被检查。用户要求“从对面看同一条街，两辆车的左右关系随视角改变”时，会暴露两个缺口：字面请求需要先推理成明确的视角与位置条件，而初始图片即使依据正确条件生成，仍可能把两车画在错误位置。前一个问题发生在生成之前，后一个问题只有看见已生成图片之后才能判断。2026年预印本UniReason 1.0的主要想法是把这两步都交给同一个统一模型：生成前先写出知识推理，把隐含要求变成明确的生成条件；生成后再读取自己的草图，写出偏差说明，并据此做一轮编辑修订。两种行为都用构造的交错样本共同训练。所核查版本为v1，论文实现的是一轮修订，不推广成任意次数的自我改进\scite{vision-unireason-2026}。

模型基于已有的Bagel检查点，而不是给Show-o2追加一个反思头。Bagel采用Transformer混合结构（mixture of transformers, MoT）：理解与生成两组专家参数在共同的交错上下文中计算，理解路径读取ViT编码的视觉特征并预测文字，生成路径在VAE潜空间预测流匹配速度。这里的“专家”指面向不同输入与目标的参数路径，不是外部审稿者。与Show-o2相比，两条路径的视觉输入不同（ViT特征与VAE潜表示），参数也不全共享，但所有专家在同一注意力上下文中互相读取。这一结构使生成出来的草图可以被重新编码为理解路径的输入，于是反思所需的“读取自己的输出”成为同一模型内的一步计算。UniReason论文没有重新报告Bagel的参数量、层数和宽度，只说明理解路径使用ViT编码器、生成路径在VAE潜空间中工作；这些结构数字须查Bagel自身的报告，本节不补。表\ref{tab:vision-models-unireason-config}按论文公开的信息汇总组件与训练安排。

\begin{table}[htbp]
  \centering
  \small
  \begin{tabularx}{\linewidth}{@{}p{14mm}XXp{18mm}@{}}
    \toprule
    组件或阶段 & 输入与输出 & 结构与规模 & 训练状态\\
    \midrule
    理解专家 & ViT视觉特征与文字$\to$推理与偏差说明 & Bagel的理解路径；规模论文未给出 & 第一阶段冻结，第二阶段训练\\
    生成专家 & VAE潜表示与上下文$\to$流匹配速度 & Bagel的生成路径；规模论文未给出 & 两阶段训练\\
    VAE & 图像$\leftrightarrow$潜表示 & 沿用Bagel & 冻结\\
    第一阶段 & 约700万文生图与50万编辑样本 & 3万次迭代，最大学习率$5\times10^{-5}$ & 只训生成专家\\
    第二阶段 & 15万知识推理文生图、10万推理编辑、3.6万交错文生图、1万交错编辑样本 & 1万次迭代，最大学习率$2\times10^{-5}$，打包序列长5万词元 & 全部专家\\
    \bottomrule
  \end{tabularx}
  \caption{UniReason 1.0的组件与训练安排，依据论文v1。第二阶段数据合计约30万条；论文实现只做一轮修订。}
  \label{tab:vision-models-unireason-config}
\end{table}

训练数据有两类样本。知识推理样本包含请求$c$、目标文字推理$T$与目标图像$\bm I$，编辑样本还包含原图$\bm I_{\mathrm{src}}$；例如请求涉及相反视角时，$T$把场景对象、观察方向和左右关系写成可用的生成条件。原论文的文生图数据从人工选择的知识种子开始，由Gemini-2.5 Pro扩充请求并写出推理，再由Qwen-Image形成图片；编辑数据以UniREdit-Data-100K为基础补充推理。反思样本则记录一次修订过程：初始生成器给出$T^{(0)}$与草图$\bm I^{(0)}$；Gemini-2.5 Pro比较图片与请求，核验对象是否存在、属性是否正确以及风格和审美偏差，写出可执行的偏差说明$T^{(1)}$；Qwen-Image-Edit据此形成$\bm I^{(1)}$，再经比较判断只保留被认为改善且符合反馈的结果。人工种子、已有记录、合成教师和筛选器承担不同职责，筛选通过仍不等于每一条合成知识都是事实真值。

训练分两阶段。第一阶段固定理解分支，只更新生成分支，用不含推理文字的已有文生图和编辑数据加强条件落实。第二阶段放开理解与生成参数，用知识推理和反思样本共同微调：知识样本同时监督$T$和$\bm I$；反思样本把$T^{(0)},\bm I^{(0)}$作为只读上下文，只监督$T^{(1)},\bm I^{(1)}$，避免把待纠正的草图也当成应模仿的输出。把文字目标位置记为$\mathcal T$、被监督图像集合记为$\mathcal G$，第$g$幅目标图片的潜向量为$\bm z_{g,1}$，噪声为$\bm z_{g,0}$，可见上下文为$C_g$；按本章噪声到数据的记号，
\begin{equation}
  \begin{aligned}
    L_{\mathrm{text}}
      &=-\sum_{j\in\mathcal T}
           \log p_{\bm\theta}(y_j\mid C_j),\\
    L_{\mathrm{img}}
      &=\sum_{g\in\mathcal G}\mathbb E_{s,\bm z_{g,0}}
        \left\lVert
          \bm v_{\bm\theta}(\bm z_g(s),s,C_g)
              -(\bm z_{g,1}-\bm z_{g,0})
        \right\rVert_2^2,\\
    L_{\mathrm{joint}}
      &=\lambda_{\mathrm{text}}L_{\mathrm{text}}
         +\lambda_{\mathrm{img}}L_{\mathrm{img}}.
  \end{aligned}
  \label{eq:vision-unireason-joint-objective}
\end{equation}
这里$C_j$只含当前目标位置之前允许读取的文字与图片，图像项使用式\eqref{eq:vision-flow-path}的中间状态。上式是按本章记号整理的监督关系，不声称论文给每个像素单独定义了知识正确性损失：文字负对数概率训练参考推理和反思措辞，图像速度误差训练目标媒体，是否符合知识及请求由数据筛选保证。原论文第二阶段取$\lambda_{\mathrm{text}}=2$、$\lambda_{\mathrm{img}}=1$，用Adam、学习率预热与余弦日程；第一阶段30K次迭代，第二阶段10K次迭代；VAE保持冻结。若反思中的某目标词元概率为0.5，其负对数概率为0.6931；被监督修订图像的一个二维潜单元残差为$(0.1,-0.2)^\top$，平方误差为0.05；仅计这两个教学单元，加权和为$2\times0.6931+0.05\approx1.4363$。它说明文字与图像误差怎样在同一次参数更新中共同起作用。

推理没有教师和目标图像。沿街景例子作一次教学推演，只考虑观察方向改变，画面横向坐标映射为$x'=w-x$，不涉及遮挡及透视。原图汽车中心横坐标为900、自行车为300，宽度$w=1200$，新视角条件应是汽车在300、自行车在900。模型先生成这段明确条件$T^{(0)}$，再从噪声经流更新得到$\bm I^{(0)}$。假设草图仍把汽车放在右侧，模型重新编码草图，写出“汽车与自行车的左右位置仍按原视角，应交换”，再把草图和这条说明送入编辑路径生成$\bm I^{(1)}$。修订是带有草图条件的另一次完整生成，需额外编码、文字预测、流求解和解码；它既可能纠正位置，也可能损坏原来正确的颜色或对象，验收应同时检查原条件保持与新增修改。真正的相反视角还有遮挡与透视变化，$x'=w-x$只用于教学复算。

\begin{table}[htbp]
  \centering
  \small
  \begin{tabularx}{\linewidth}{@{}p{18mm}XXX@{}}
    \toprule
    比较项 & 接口调用 & 共享表示（Show-o2） & 反思反馈（UniReason）\\
    \midrule
    传递的内容 & 文字说明、掩码或适配特征 & 主干中的共享视觉表示与上下文 & 共享上下文，外加草图与偏差说明\\
    共享参数 & 无，或只训练适配器 & 主干共享，输出头分开 & 专家路径分开，同一注意力上下文\\
    训练 & 各自独立训练 & 文字与速度误差共同训练 & 两阶段，加入推理与反思样本\\
    生成后检查 & 需外部步骤 & 无 & 一轮自我读取与修订\\
    替换模块 & 容易 & 需重训 & 需重训\\
    \bottomrule
  \end{tabularx}
  \caption{理解与生成三种结合方式在同一图片修改请求上的比较。}
  \label{tab:vision-models-unified-compare}
\end{table}

表\ref{tab:vision-models-unified-compare}在同一图片修改请求下比较三种结合方式。从接口调用到共享表示，改变的是理解结果的传递方式：理解看到的细节不再必须转述成文字，而在主干中直接可用，代价是两种能力共用参数、可能互相干扰，且无法单独替换生成器。从共享表示到反思反馈，改变的是信息流的方向：生成的图片回到理解路径被检查，错误可以在交付前得到一次修正，代价是每次修订都要再运行一遍编码、文字预测和生成，并依赖合成教师构造的训练数据。结合越紧，能在两类能力之间流动的信息越多，但出错时越难定位到某一个模块。

\section{模型的使用约束}
\label{sec:vision-models-constraints}

一个模型支持的输入和输出，是开始选型的依据，但不是能否完成任务的证明。实际使用还需回答三个问题：编码、压缩和生成之后，当前任务需要的信息是否还在；模型声称接受的各个模态，在一次回答中是否真的被使用；达到这一结果用了多少计算资源。三个问题分别对应信息、模态和资源三个维度，下面回到同一街景与音画请求逐一检查。

\subsection{信息保留的约束}
\label{sec:vision-models-information}

整图缩小会让车牌字符变成几个像素，局部裁剪会删除周围关系，视频抽帧会漏掉短动作，固定查询数会把许多视觉位置汇总成少量表示，潜空间压缩会丢掉生成所需的细纹理。投影和共享主干可以重新组织已经保留的信息，却不能恢复这些阶段丢弃的证据。可以固定问题，分别提供整图、局部高分辨率图，以及带原图位置的二者组合。如果只有局部图能正确读取车牌，就应进一步判断错误来自原始分辨率、输入缩放还是后续压缩，而不是立即归为语言推理失败。组合输入还要确认模型知道局部图来自哪里，否则可能把两个视图当成两辆不同的车。

输出也需逐项验收。匹配分数正确只说明当前候选中的选择，语言流畅只说明形式自然，视频画面逼真只说明感知质量的一部分；数量、位置、时间与来源需要对应的检查。精确区域和答案证据见\ref{chap:vision-multimodal-understanding}章，媒体连贯性和控制见\ref{chap:vision-multimodal-generation}章。

\subsection{模态贡献的约束}
\label{sec:vision-models-modality-contribution}

模型接受声音，并不说明一次回答确实使用了声音。可以在保持问题不变时移除音频，或换成时间错开的声音，检查结果与正确答案怎样变化。若问题只问车的颜色，移除声音后答案不变是合理的；若问题问敲击发生了几次，错误声音改变了答案，则需要判断模型是否把它与画面进行了适当核验。同样，替换街景图却保留文字问题，可以检查模型是否依赖常见语言模式猜测。替换条件必须使正确答案或证据关系发生可解释的变化；只观察回答变了，无法证明模型使用了正确模态。CounterCredit用随机图块替换真实返回，正是这种对照在训练中的用法。

模态贡献还受训练依据限制。整图配对监督适合建立全局对应，框或时间区间标注提供更细的位置联系，指令示例说明请求和回答怎样组织；它们可以联合使用，但不能互相冒充。ImageBind中未直接配对模态之间的联系、Show-o2的视频能力，都只在相应训练覆盖的范围内成立。多模态评价需要覆盖缺失模态、冲突模态和训练中未直接出现的组合，不能由一次演示推断任意模态都已对齐。

\subsection{计算资源的约束}
\label{sec:vision-models-cost}

多模态成本可以按输入编码、联合上下文计算、输出生成和媒体解码记录。固定图像大小还不够，裁剪数、视频采样、音频长度、历史轮数和输出时长都可能改变实际输入与输出规模。一个只生成短答案的调用，不能与生成数秒音视频的调用仅比较总延迟。在标准稠密注意力中，$n$个序列位置的成对分数规模为$n^2$，隐藏宽度为$d$时，相关乘加量含$O(n^2d)$项。若有256个文字单元和1024个视觉单元，总长为1280；视觉单元增加到4096时，总长4352，成对位置数约变为原来的11.56倍。这正是BLIP-2和Flamingo固定查询数、Qwen3-VL合并相邻特征、Wan在时间上压缩的原因；这个比例只描述该计算项，局部注意力、缓存与分块实现都会改变实际延迟。

训练成本与推理成本也要分开。冻结模块减少需要保存的参数梯度和优化器状态，却不省去推理中的前向计算；若可训练连接器位于冻结语言模型之前，训练还要经过冻结层计算输入梯度。媒体生成反复调用速度或去噪网络，成本随步数和每步条件计算增长，无分类器引导使每步前向加倍，候选重排则重复整条生成路径。结合方式也有各自成本：共享主干可以复用上下文，接口调用可以独立缓存编码或选择不同规模的专用模块，反思修订则把生成成本至少加倍。部署时应记录编码、缓存和生成各环节的资源消耗，以及等待和输出长度的限制。

\section{本章小结}
\label{sec:vision-models-summary}

开篇的问题是文字、图像和声音怎样进入同一个模型、在哪里发生联系、又怎样恢复成结果。编码给出第一部分答案：文字切成词元，声音转为梅尔频谱，图像和视频沿用上一章的切块编码，每个单元都要保留来源、位置和时间。连接给出第二部分答案：独立编码后在共同空间比较，让语言计算读取视觉特征，或在共享主干中交互，三种安排在可预计算性和逐位置读取能力之间取舍。解码给出第三部分答案：分数、词元和媒体潜表示各需相应的输出头与解码器。

按输出类型，本章比较了三组模型。CLIP与ImageBind用对比目标形成可比较表示，后者以图像为中心把更多模态接进已有空间。语言响应模型都复用已有视觉编码器与语言模型，差别在视觉信息进入语言计算的位置和数量，CounterCredit进一步让模型在回答中途决定是否取回新像素，并用决策和证据两个对照约束这种观察。图像生成模型分别改变生成空间（LDM）、中间表示（unCLIP）和采样引导（GLIDE）；视频与音频生成以流匹配为共同目标，Wan用掩码统一多种视频条件，LTX-2用双分支让画面和声音共同形成。理解与生成可以通过接口调用、共享表示或反思反馈结合，结合越紧，传递的信息越多，模块也越难单独替换和诊断。各模型的配置表还显示了一个共同规律：新训练的连接部分通常远小于被复用的主干，LLaVA只有一个投影矩阵，BLIP-2的Q-Former约1.88亿参数，Flamingo的门控层随语言模型从12亿增到100亿，而冻结或沿用的视觉编码器与语言模型占去其余大部分参数；连接的设计因此决定了少量可训练参数能把多少视觉证据交给主干。

还有一些问题本章没有解决。模型的输入范围不等于证据被保留，接受某种模态不等于使用了它，结合理解与生成也不会自动消除证据损失与输出错误；怎样在固定预算中保留必需的局部证据、建立可靠的时间对应，并协调不同输出目标，仍需结合具体任务研究。后续任务章在这些模型接口上加入各自的观测、控制、标注与成本要求，并检查答案与证据、内容与条件是否一致。

\paragraph{扩展阅读}
图文表示方面，ALIGN与CLIP同期，它不对网页替代文字做精细清洗，用约18亿对含噪图文对训练双塔对比模型，说明数据规模可以在一定程度上抵消配对噪声；Flamingo训练数据中的ALIGN图文对即来自这项工作\scite{vision-2102-05918}。SigLIP把CLIP沿行、列归一化的对比损失换成对每一对图文独立判断是否匹配的sigmoid损失，不再需要在整个相似度矩阵上归一化，在较小批量下效果更好；Show-o2复用的视觉层即来自SigLIP\scite{vision-2303-15343}。CoCa在对比目标之外加入以交叉注意力读取图像的描述生成解码器，在一次预训练中同时得到可比较的整图表示和生成描述的能力\scite{vision-2205-01917}。

语言响应方面，PaLI联合扩大约40亿参数的ViT视觉编码器与mT5语言模型，并在覆盖百余种语言的WebLI图文数据上训练，表明视觉侧的规模同样影响描述和问答\scite{vision-2209-06794}。InternVL的出发点是视觉编码器的规模远落后于语言模型，它把视觉编码器扩大到60亿参数，先以对比学习与语言侧对齐、再以生成目标对齐，然后接入大语言模型；这与本章BLIP-2、LLaVA等只训练小型连接部分的做法形成对照\scite{vision-2312-14238}。

统一与生成方面，Chameleon把图像也离散为码，与文字词元混排在同一序列中，由单个自回归Transformer从头训练，代表在输入端即融合的路线\scite{vision-2405-09818}。Transfusion在同一个Transformer中对文字做下一词元预测、对图像的连续潜表示做扩散，与Show-o2同属基于共享表示的结合，但不要求把图像离散化\scite{vision-2408-11039}。Imagen在像素空间做级联扩散，发现扩大冻结的纯文本编码器T5比扩大图像扩散模型更能提高图文一致性，与unCLIP依赖CLIP表示的选择形成对照\scite{vision-2205-11487}。Parti则沿离散码自回归路线，把文字到图像视为序列到序列的翻译，并把模型扩大到200亿参数\scite{vision-2206-10789}。希望系统了解这一领域，可阅读Li等的综述，它按视觉理解、视觉生成、统一视觉模型、借助大语言模型的多模态模型和多模态智能体梳理了相关工作\scite{vision-2309-10020}。
